% Report tecnico-scientifico sull'archivio GRIB2 CHAPTER.
% Compilazione: ./compila.sh (latexmk + biber da TinyTeX in ~/.TinyTeX).
\documentclass[11pt,a4paper]{report}

\usepackage[utf8]{inputenc}
\usepackage[T1]{fontenc}
\usepackage{lmodern}
\usepackage[italian]{babel}
\usepackage[autostyle]{csquotes}
\usepackage{microtype}
\usepackage[margin=2.5cm]{geometry}

\usepackage{amsmath}
\usepackage{siunitx}
\sisetup{locale=IT, group-minimum-digits=4, per-mode=symbol, exponent-product=\cdot, range-phrase={--}, range-units=single}
\DeclareSIUnit{\hPa}{hPa}

\usepackage{booktabs}
\usepackage{array}
\usepackage{tabularx}
\usepackage{longtable}
\usepackage{caption}
\captionsetup{font=small, labelfont=sc}

\setlength{\marginparwidth}{2cm}
\usepackage[textsize=footnotesize, colorinlistoftodos]{todonotes}

\usepackage[backend=biber, style=numeric-comp, sorting=none, maxbibnames=10]{biblatex}
\addbibresource{bib/riferimenti.bib}

\usepackage[hidelinks]{hyperref}
\usepackage[capitalise]{cleveref}

% Nomi di campi, variabili e chiavi GRIB, scritti come nel file.
\newcommand{\campo}[1]{\texttt{\detokenize{#1}}}

\title{L'archivio GRIB2 CHAPTER\\[0.5em]
  \large Downscaling dinamico a \SI{3}{km} di ERA5\\[0.3em]
  Contenuto e corrispondenza con lo schema ERA5}
\author{Luca Monaco\\Fondazione CIMA}
\date{ottobre 2026}

\begin{document}

\pagenumbering{alph}
\maketitle
\pagenumbering{roman}
\tableofcontents
\cleardoublepage
\pagenumbering{arabic}

\chapter{La simulazione}
\label{cap:simulazione}

CHAPTER è un downscaling dinamico a convezione esplicita di ERA5 sull'Europa e sul bacino del
Mediterraneo \parencite{bernini2025}, prodotto con il modello WRF (Weather Research and
Forecasting), versione 4.1.1 \parencite{wrf411}. L'archivio GRIB2 CHAPTER è l'insieme dei file
orari GRIB2 ricavati dai file wrfout della simulazione, uno per istante di validità, con i nomi
(\campo{shortName}) e i codici (\campo{paramId}) dei parametri ECMWF (\cref{cap:formato}).

\section{Dominio, griglia e periodo}
\label{sec:dominio}

La simulazione usa il dominio d02, una griglia Mercator di $\num{1641} \times \num{1353}$ celle
con passo di \SI{3}{km}, cioè \num{2220273} punti per campo. Il parallelo standard della
proiezione è \ang{44.671} N. Il dominio va da \ang{23.4719} a \ang{60.0006} di latitudine nord e
da \ang{19.9722} O a \ang{42.2555} E. In verticale il modello ha 49 livelli, con il top a
\SI{50}{\hPa}. WRF integra le equazioni su una sfera di raggio \SI{6370}{km}, e l'archivio dichiara
la stessa sfera nella definizione di griglia (\cref{sec:geometria}).

L'archivio copre a passo orario gli anni 2019 e 2024 e il periodo dal 1º gennaio al 2 luglio 2025,
per \num{21936} istanti di validità. Il 2 luglio 2025 alle 23 UTC è l'ultimo istante prodotto dalla
simulazione. Ogni run giornaliera parte alle 18 UTC del giorno che precede quello a cui è dedicata
e integra \SI{30}{h}; le prime \SI{6}{h} sono di spin-up e non entrano nell'archivio. Il giorno D è
fatto degli istanti di validità dalle 01 alle 23 UTC della run partita alle 18 UTC di D$-1$, e il suo
istante di validità delle 00 UTC è l'ultima ora, la trentesima, della run partita alle 18 UTC di
D$-2$ (\cref{sec:tempo,sec:cambio-run}).

\section{Configurazione fisica}
\label{sec:configurazione}

La \cref{tab:configurazione} riporta le opzioni della simulazione da cui dipende il contenuto
dell'archivio, lette dagli attributi globali dei file wrfout. Sono le stesse nei tre periodi
controllati, il 2019, marzo 2024 e luglio 2024.

\begin{table}[htbp]
  \centering
  \small
  \caption{Opzioni della simulazione che determinano quali campi l'archivio può contenere.}
  \label{tab:configurazione}
  \begin{tabularx}{\textwidth}{@{}lllX@{}}
    \toprule
    Opzione & Valore & Schema & Conseguenza per l'archivio \\
    \midrule
    \campo{BL_PBL_PHYSICS}     & 1 & YSU      & schema non locale, senza energia cinetica turbolenta prognostica \\
    \campo{KM_OPT}             & 4 &          & la diffusione non produce energia cinetica turbolenta \\
    \campo{SF_OCEAN_PHYSICS}   & 0 &          & oceano non simulato: niente correnti né altezza della superficie del mare \\
    nessun modello d'onda      &   &          & nessun campo d'onda \\
    \campo{SST_UPDATE}         & 0 &          & temperatura del mare costante in ogni run giornaliera \\
    \campo{CU_PHYSICS}         & 0 &          & convezione risolta; \campo{RAINC} identicamente nullo \\
    \campo{MP_PHYSICS}         & 6 & WSM6     & acqua di nube, pioggia, ghiaccio, neve e graupel; nessuna grandine \\
    \campo{ICLOUD}             & 1 & Xu--Randall & frazione nuvolosa continua fra 0 e 1 \\
    \campo{SF_SURFACE_PHYSICS} & 3 & RUC      & suolo su sei nodi a \SIlist{0;5;20;40;160;300}{cm} \\
    \campo{SF_LAKE_PHYSICS}    & 0 &          & nessun modello di lago \\
    \campo{SF_URBAN_PHYSICS}   & 0 &          & nessuno schema urbano \\
    \campo{aer_opt}            & 0 &          & radiazione senza aerosol \\
    \campo{usemonalb}          & \campo{.false.} & & albedo di fondo da tabella, non da satellite \\
    \campo{rdlai2d}            & \campo{.false.} & & indice di area fogliare da tabella per categoria \\
    \campo{isfflx}             & 1 &          & flussi turbolenti di superficie calcolati a ogni passo temporale \\
    \campo{output_diagnostics} & 0 &          & nessun estremo orario di temperatura a \SI{2}{m} \\
    \campo{nwp_diagnostics}    & 1 &          & massimo orario del vento a \SI{10}{m} e altri sei massimi orari \\
    \bottomrule
  \end{tabularx}
\end{table}

Lo schema radiativo è RRTMG, con la climatologia dell'ozono di CAM (\campo{o3input=2}). La
microfisica WSM6 porta cinque specie condensate, fra cui il graupel, che ERA5 non pubblica come
campo separato (\cref{sec:colonne}).

Nello schema di superficie RUC l'albedo di fondo, l'indice di area fogliare e la rugosità sono
valori di tabella per categoria dominante, cioè per la classe di uso del suolo con la frazione
maggiore nella cella. Per questo \campo{al} prende da 8 a 11 valori su tutto il dominio e
\campo{fsr}, fuori dalla neve, un valore per categoria (\cref{cap:superficie,cap:suolo}).

L'energia cinetica turbolenta, le onde, le correnti e l'altezza della superficie del mare non
esistono in questa simulazione. L'altezza dello strato limite e i flussi di calore sensibile e
latente, invece, sono stati calcolati ma non scritti nei file wrfout (\cref{cap:limiti}).

\section{Il file statico}
\label{sec:file-statico}

Le informazioni statiche su suolo, vegetazione e laghi vengono dal file statico
\campo{geo_em.d02}, prodotto dal preprocessore di WRF (WPS) per questa simulazione. La sua
categoria di uso del suolo (\campo{LU_INDEX}) coincide con la categoria dominante scritta ogni ora
nei file wrfout (\campo{IVGTYP}) sul \SI{100.00}{\percent} delle celle a luglio 2019 e a gennaio
2025, una volta esclusa la riclassificazione dei laghi come acqua, che il modello applica perché
la fisica dei laghi è spenta. Il 15 febbraio 2025 alle 12 UTC l'accordo scende al
\SI{99.98}{\percent}: le \num{2735} celle diverse portano tutte ghiaccio marino, che il modello
riclassifica come terra a ogni ora (\cref{sec:statici-orari}).

L'uso del suolo è espresso in 21 classi MODIS-IGBP (\campo{MMINLU=MODIFIED_IGBP_MODIS_NOAH}). Il
file statico porta per ogni cella la frazione di ciascuna classe, mentre i file wrfout ne
conservano solo la categoria dominante. In media la categoria dominante occupa il
\SI{90.1}{\percent} di una cella di terra, e il \SI{17.93}{\percent} delle celle di terra ha sia
vegetazione alta sia vegetazione bassa sopra il \SI{5}{\percent}. Dal file statico vengono anche la
tessitura del suolo e la profondità dei laghi. Il file statico ha reso possibili quindici campi:
le coperture e i tipi di vegetazione (\campo{cvl}, \campo{cvh}, \campo{tvl}, \campo{tvh}), il tipo
di suolo (\campo{slt}), la copertura e la profondità dei laghi (\campo{cl}, \campo{dl}) e gli otto
campi di umidità e temperatura del suolo sugli strati ERA5, a cui fornisce la maschera di terra
permanente (\cref{cap:suolo}).

\chapter{Il formato dell'archivio}
\label{cap:formato}

\section{File e messaggi}
\label{sec:file}

L'archivio è fatto di un file GRIB2 per ogni istante di validità. Ogni file orario contiene 90
variabili in 246 messaggi GRIB, sempre gli stessi: l'elenco viene da un unico registro del
convertitore, e ogni file prodotto viene confrontato con quell'elenco, così un campo che non si è
potuto calcolare non passa inosservato. I nomi dei file seguono lo schema
\campo{ailam-an-cima-3km-{year}-{year}-1h-v1-{YYYYMMDD}{HH}.grib}, in cui l'anno compare due volte
con lo stesso valore e la parte finale è l'istante di validità.

I messaggi usano il template di griglia 3.10 (Mercator), il centro di origine \campo{ecmf} e la
compressione \campo{grid_ccsds} a 24 bit per valore. Ogni parametro porta il paramId e lo
shortName ECMWF \parencite{paramdb}. Gli shortName distinti sono 89 per 90 variabili, perché
\campo{z} compare due volte, come geopotenziale sui livelli di pressione e come geopotenziale di
superficie, cioè l'orografia. I campi definiti solo su una parte del dominio portano una bitmap
GRIB, e i punti fuori dal dominio di definizione sono punti mancanti, non zeri
(\cref{sec:mascheratura}).

Un file orario pesa in media \SI{737}{MB} sui 184 file di settembre 2024, e \SI{661}{MB} sugli 8
file delle 00 UTC fra questi; sui 476 file di febbraio e marzo 2024 la media è \SI{765}{MB}. Ne
risultano fra \SI{17.7}{GB} e \SI{18.4}{GB} al giorno e fra \SI{6.5}{TB} e \SI{6.7}{TB} all'anno.

\section{Geometria}
\label{sec:geometria}

La griglia dichiara la forma della Terra con \campo{shapeOfTheEarth=1} e raggio
\SI{6370000}{m}, la sfera del modello (\cref{sec:dominio}). Con questa dichiarazione le coordinate
che eccodes ricostruisce dalla definizione di griglia coincidono con quelle del modello entro
\SI{2.009}{m} su tutto il dominio. Senza la dichiarazione, la griglia ricostruita da eccodes si
sposta fino a circa \SI{1.1}{km} al bordo nord.

\section{Livelli verticali}
\label{sec:livelli}

Le 13 variabili tridimensionali stanno sugli stessi 13 livelli di pressione,
\SIlist{1000;925;850;700;600;500;400;300;250;200;150;100;50}{\hPa}. Sono i livelli richiesti dal
progetto LAM AI, un sottoinsieme dei 37 livelli di pressione di ERA5 \parencite{era5doc}. Degli
altri, gli 8 sopra \SI{50}{\hPa} stanno sopra il top del modello e non si possono produrre.

I 246 messaggi si ripartiscono fra i tipi di livello della \cref{tab:tipi-livello}.

\begin{table}[htbp]
  \centering
  \small
  \caption{Tipi di livello dei 246 messaggi di un file orario.}
  \label{tab:tipi-livello}
  \begin{tabularx}{\textwidth}{@{}lrX@{}}
    \toprule
    Tipo di livello & Messaggi & Contenuto \\
    \midrule
    \campo{isobaricInhPa}      & 169 & 13 variabili su 13 livelli di pressione \\
    \campo{heightAboveGround}  & 11  & vento a 10, 100 e \SI{200}{m}; temperatura, temperatura di rugiada e umidità relativa a \SI{2}{m}; \campo{10fg}; \campo{vwsh} \\
    \campo{surface}            & 47  & campi statici, suolo, neve, precipitazione, deflusso e radiazione in superficie \\
    \campo{entireAtmosphere}   & 10  & integrali di colonna e coperture nuvolose \\
    \campo{nominalTop}         & 5   & radiazione al top dell'atmosfera \\
    \campo{mostUnstableParcel} & 2   & \campo{mucape}, \campo{mucin} \\
    \campo{meanSea}            & 1   & \campo{msl} \\
    \campo{entireLake}         & 1   & \campo{dl} \\
    \bottomrule
  \end{tabularx}
\end{table}

\section{Convenzione temporale}
\label{sec:tempo}

I campi istantanei usano il template di prodotto 4.0, e il loro tempo di riferimento
(\campo{dataDate}, \campo{dataTime}) coincide con l'istante di validità. I 19 campi accumulati
sono la precipitazione totale (\campo{tp}), la pioggia (\campo{tirf}), la neve
(\campo{sf}), i tre deflussi (\campo{sro}, \campo{ssro}, \campo{ro}) e i 13 flussi radiativi in
superficie e al top dell'atmosfera. Usano il template 4.8, perché le definizioni ECMWF di questi
parametri richiedono la chiave \campo{typeOfStatisticalProcessing}, che nel template 4.0 non esiste
\parencite{eccodes}. Come in ERA5 \parencite{era5doc}, ogni campo accumulato copre l'ora
che termina all'istante di validità $H$; nell'archivio il tempo di riferimento è $H-1$ e
l'intervallo di accumulazione (\campo{stepRange}) è 0--1. Il valore è la differenza fra due ore successive della stessa run, per
esempio $\campo{tp}(H) = [\campo{RAINNC}(H) - \campo{RAINNC}(H-1)]/1000$, con \campo{RAINNC} in
\si{mm} e \campo{tp} in \si{m} \parencite{issue70}. Il ventesimo campo con template 4.8 è
\campo{10fg}, il massimo del vento a \SI{10}{m}, con tempo di riferimento $H-1$, intervallo
dichiarato 0--1 e \campo{stepType=max} (\cref{cap:raffica}).

Il file orario delle 00 UTC del giorno D viene per intero, campi accumulati compresi, dalla
trentesima ora della run partita alle 18 UTC di D$-2$; i campi accumulati sono la differenza fra la
trentesima e la ventinovesima ora \parencite{issue71,issue72}. I campi accumulati del file orario
delle 01 UTC di D sono la differenza fra la settima e la sesta ora della run partita alle 18 UTC di
D$-1$. I 24 istanti di validità dalle 01 UTC di D alle 00 UTC di D$+1$ vengono tutti da questa run, e
la somma dei loro campi accumulati è il totale della run dalla sesta alla trentesima ora. Quando il
file wrfout della trentesima ora di una run manca alla sorgente, ogni messaggio del file orario
delle 00 UTC successivo è fatto di punti mancanti; al 14 ottobre 2026 non succede in nessun file
orario dell'archivio.

Poiché i campi istantanei hanno tempo di riferimento $H$ e i campi con template 4.8 hanno $H-1$, un
file orario contiene due tempi di riferimento, e l'indice dei file orari va costruito sull'istante
di validità (\campo{validityDate}, \campo{validityTime}), mai sul tempo di riferimento.
Per distinguere un campo accumulato da uno istantaneo servono \campo{stepType} (o
\campo{typeOfStatisticalProcessing}) e \campo{stepRange}. La chiave
\campo{generatingProcessIdentifier} vale 128, il valore predefinito di eccodes per GRIB2, in tutti
i 246 messaggi, e non separa i campi accumulati da quelli istantanei \parencite{issue53}.

\section{Il cambio di run}
\label{sec:cambio-run}

Il \emph{cambio di run} è il passaggio, una volta al giorno, fra l'istante di validità delle 00 UTC
di D, trentesima ora della run partita alle 18 UTC di D$-2$, e quello delle 01 UTC di D, settima ora
della run partita alle 18 UTC di D$-1$. Il salto fra i due istanti di validità è la differenza fra
le due run alle 00 UTC di D, trentesima ora della prima e sesta della seconda, più la variazione
dell'ora 00--01 UTC nella seconda. La \cref{tab:cambio-run} lo
confronta con la variazione di un'ora dentro la stessa run, sul dominio intero, su un solo cambio di
run, dalle 23 UTC del 18 gennaio alle 01 UTC del 19 gennaio 2024 \parencite{issue72}.

\begin{table}[htbp]
  \centering
  \small
  \caption{Scarto quadratico medio fra due istanti di validità consecutivi, sul dominio intero, 18 e
  19 gennaio 2024.}
  \label{tab:cambio-run}
  \begin{tabular}{@{}lSSSSS@{}}
    \toprule
    Istanti & {\campo{2t} (\si{K})} & {\campo{sp} (\si{\hPa})} & {modulo del vento a \SI{10}{m} (\si{m/s})} & {\campo{t} a \SI{850}{\hPa} (\si{K})} & {\campo{tcwv} (\si{kg/m^2})} \\
    \midrule
    23--00 UTC, stessa run & 0.66 & 0.50 & 1.13 & 0.48 & 0.89 \\
    00--01 UTC, cambio di run & 1.79 & 1.62 & 1.62 & 0.94 & 1.54 \\
    \bottomrule
  \end{tabular}
\end{table}

Il 19 gennaio 2024 lo scarto quadratico medio di \campo{2t} al cambio di run è \num{2.7} volte
quello di un'ora dentro la stessa run, e la differenza assoluta massima fra i due istanti di
validità è \SI{11.0}{K}, contro \SI{9.1}{K} fra le 23 e le 00 UTC. Una
tendenza o una derivata temporale calcolata attraverso l'intervallo fra le 00 e le 01 UTC contiene
questo salto.

\chapter{Regole di corrispondenza con ERA5}
\label{cap:regole}

Due regole decidono che cosa l'archivio pubblica e dove, e valgono per tutte le 90 variabili.

\section{Regola di pubblicazione}
\label{sec:pubblicazione}

Un campo si pubblica sotto un nome ERA5 solo se la grandezza WRF è quella che il parametro ERA5
definisce, oppure se la differenza si può dire in una riga. Un campo che non soddisfa nessuna
delle due condizioni resta fuori dall'archivio, anche quando il modello lo calcola. I campi esclusi
restano fuori per uno di due motivi, e non sono la stessa cosa: per alcuni campi WRF nessun
parametro ECMWF descrive la stessa grandezza, per altri il parametro esiste ma ERA5 non lo pubblica
(\cref{sec:non-pubblicati}).

Sette campi, in sei righe della \cref{tab:approssimazioni}, sono pubblicati con una differenza
dichiarata rispetto al loro omonimo ERA5.

\begin{table}[htbp]
  \centering
  \small
  \caption{Campi pubblicati con una differenza dichiarata rispetto al parametro ERA5.}
  \label{tab:approssimazioni}
  \begin{tabularx}{\textwidth}{@{}lXXl@{}}
    \toprule
    Campo & Parametro ERA5 & Campo CHAPTER & Sezione \\
    \midrule
    \campo{slor} & pendenza dell'orografia sottogriglia & pendenza del terreno risolto a \SI{3}{km} & \ref{sec:orografia} \\
    \campo{10fg} & raffica parametrizzata & massimo orario del vento risolto a \SI{10}{m} & \ref{cap:raffica} \\
    \campo{iews}, \campo{inss} & componenti dello sforzo turbolento in superficie dell'IFS & $\rho u_*^2$, con $\rho$ densità dell'aria e $u_*$ velocità d'attrito, proiettato sulla direzione del vento a \SI{10}{m} & \ref{sec:rugosita} \\
    \campo{fsr} & rugosità efficace, con la resistenza orografica parametrizzata & rugosità locale \campo{ZNT}, valore di tabella della categoria dominante modificato dalla neve & \ref{sec:rugosita} \\
    \campo{lsm} & maschera terra--mare permanente & maschera oraria di WRF, che conta come terra il ghiaccio marino & \ref{sec:statici-orari} \\
    \campo{tcw} & acqua totale in colonna, cinque specie & somma di sei specie, graupel incluso & \ref{sec:colonne} \\
    \bottomrule
  \end{tabularx}
\end{table}

\section{Regola di mascheratura}
\label{sec:mascheratura}

Un punto è mancante dove il parametro ERA5 non è definito. Un punto in cui il modello non fornisce
un valore, o in cui il valore sembra sbagliato, non per questo è mancante. I punti mancanti sono
scritti come bitmap GRIB, mai come zeri, e chi legge l'archivio distingue sempre un punto in cui il
parametro non è definito da uno in cui vale zero. La \cref{tab:maschere} elenca i campi mascherati.

\begin{table}[htbp]
  \centering
  \small
  \caption{Campi con punti mancanti e dominio su cui sono definiti.}
  \label{tab:maschere}
  \begin{tabularx}{\textwidth}{@{}lXl@{}}
    \toprule
    Campo & Definito su & Sezione \\
    \midrule
    \campo{sst}, \campo{ci} & celle di mare & \ref{sec:mare} \\
    \campo{dl} & celle con copertura di lago maggiore di zero (\num{32000}) & \ref{sec:vegetazione} \\
    \campo{swvl1-4}, \campo{stl1-4}, \campo{slt} & terra permanente (\num{1304109} celle) & \ref{sec:suolo}, \ref{sec:vegetazione} \\
    \campo{tvl} & tutte le celle, tranne le \num{124672} in cui domina un arbusteto & \ref{sec:vegetazione} \\
    \campo{mucin} & colonne con CAPE di almeno \SI{100}{J/kg} & \ref{sec:instabilita} \\
    \campo{tsn} & celle con \campo{SNOWC} sopra \num{0.9} & \ref{sec:neve} \\
    \campo{fal} & punti con radiazione solare incidente in superficie sopra \SI{50}{W/m^2} & \ref{sec:albedo} \\
    \bottomrule
  \end{tabularx}
\end{table}

La copertura dei laghi \campo{cl} è definita anche sul mare, dove vale zero, perché una cella di
mare non ha copertura lacustre. Quattro campi sono mascherati per una ragione diversa: nei punti mancanti
la grandezza WRF non ha un valore che corrisponda al parametro ERA5. \campo{tvl} manca sugli
arbusteti, perché la tabella dei tipi ERA5 separa arbusti sempreverdi e decidui e la classificazione
MODIS non porta la fenologia degli arbusti (\cref{sec:vegetazione}). \campo{tsn} manca dove la neve
non copre la cella, perché lì il campo del modello è una temperatura del suolo (\cref{sec:neve}).
\campo{fal} manca dove la radiazione solare incidente in superficie non supera \SI{50}{W/m^2};
senza radiazione incidente il rapporto fra radiazione riflessa e incidente non è definito (\cref{sec:albedo}). \campo{mucin} manca dove CAPE è sotto \SI{100}{J/kg}, perché lì la routine non restituisce
un'inibizione (\cref{sec:instabilita}).

\subsection{Mare e acque interne}
\label{sec:mare}

Per mascherare \campo{sst} e \campo{ci} serve distinguere il mare dalle acque interne, e la classe
di lago del file statico non basta. La classe 21 (lago) della classificazione MODIS-IGBP comprende
il Mar Nero, il Mare d'Azov e parte del Baltico, che insieme sono l'\SI{89}{\percent} delle
\num{55881} celle che il file statico classifica interamente come lago. Mascherare su quella classe
cancellerebbe \campo{sst} dal Mar Nero.

Il criterio usato è la dimensione della componente connessa, cioè dell'insieme di celle d'acqua
contigue per lato. Nel dominio le componenti più grandi hanno \num{844431} celle (Atlantico,
Mediterraneo, Baltico e Mare del Nord), \num{49841} (Mar Nero e Mare d'Azov) e \num{8583} (Mar
Rosso). La successiva è il lago Vänern, con \num{1225} celle, il lago più grande del dominio. Sono
mare le componenti di almeno \num{2000} celle, per un totale di \num{902855} celle di mare.
Qualunque soglia fra \num{1226} e \num{8583} celle dà lo stesso insieme.

\section{Campi statici e campi orari}
\label{sec:statici-orari}

Un campo statico viene dal file statico ed è quindi identico in ogni file orario. Sono statici
\campo{cl}, \campo{cvh}, \campo{cvl}, \campo{dl}, \campo{slt}, \campo{tvh} e \campo{tvl}. Le
colonne di suolo \campo{swvl1-4} e \campo{stl1-4} cambiano ogni ora, ma la loro maschera è la terra
permanente del file statico.

La maschera terra--mare \campo{lsm} è invece la maschera oraria di WRF (\campo{LANDMASK}), che
riclassifica come terra ogni cella d'acqua coperta da ghiaccio marino. La maschera cresce quindi in
inverno: in un file del 20 marzo 2024 le celle riclassificate sono \num{2469}, esattamente le celle
con ghiaccio marino. Il ghiaccio marino \campo{ci} vale 0 o 1, perché la simulazione non usa
ghiaccio frazionario (\campo{FRACTIONAL_SEAICE=0}). La maschera di terra permanente, nel senso di
ERA5, è l'insieme dei punti con $\campo{lsm} = 1$ e $\campo{ci} = 0$, contando come zero i punti
mancanti di \campo{ci}.

Anche l'albedo di fondo \campo{al} cambia nel tempo, con il ghiaccio marino e al passaggio fra la
stagione estiva e quella invernale della tabella di uso del suolo (\cref{sec:albedo}). Per
\campo{lsm} e \campo{al} vale il file orario dell'istante di validità, non un file qualunque
dell'archivio.

\chapter{Atmosfera}
\label{cap:atmosfera}

\section{Le variabili sui livelli di pressione}
\label{sec:livelli-pressione}

La \cref{tab:var-pressione} elenca le 13 variabili sui livelli di pressione, 169 messaggi in
tutto. Oltre a temperatura, vento, umidità e geopotenziale comprendono le due forme della
velocità verticale, in pressione (\campo{w}) e geometrica (\campo{wz}), la frazione nuvolosa e
quattro contenuti d'acqua di nube e di precipitazione.

\begin{table}[htbp]
  \centering
  \small
  \caption{Variabili sui 13 livelli di pressione.}
  \label{tab:var-pressione}
  \begin{tabular}{@{}lrll@{}}
    \toprule
    Campo & paramId & Descrizione & Unità \\
    \midrule
    \campo{cc}   & 248    & frazione nuvolosa & 0--1 \\
    \campo{ciwc} & 247    & contenuto specifico di ghiaccio di nube & \si{kg/kg} \\
    \campo{clwc} & 246    & contenuto specifico di acqua liquida di nube & \si{kg/kg} \\
    \campo{crwc} & 75     & contenuto specifico di pioggia & \si{kg/kg} \\
    \campo{cswc} & 76     & contenuto specifico di neve & \si{kg/kg} \\
    \campo{q}    & 133    & umidità specifica & \si{kg/kg} \\
    \campo{r}    & 157    & umidità relativa & \si{\percent} \\
    \campo{t}    & 130    & temperatura & \si{K} \\
    \campo{u}    & 131    & componente zonale del vento & \si{m/s} \\
    \campo{v}    & 132    & componente meridiana del vento & \si{m/s} \\
    \campo{w}    & 135    & velocità verticale in pressione & \si{Pa/s} \\
    \campo{wz}   & 260238 & velocità verticale geometrica & \si{m/s} \\
    \campo{z}    & 129    & geopotenziale & \si{m^2/s^2} \\
    \bottomrule
  \end{tabular}
\end{table}

\paragraph{Contenuti specifici.} WRF porta vapore e idrometeore come \emph{rapporti di mescolanza},
per chilogrammo di aria secca. ECMWF definisce \campo{q} e i contenuti d'acqua come \emph{contenuti
specifici}, per chilogrammo di aria umida, e nella massa di aria umida conta aria secca, vapore, acqua di nube,
ghiaccio, pioggia e neve \parencite{paramdb}. L'archivio divide quindi \campo{q}, \campo{clwc},
\campo{ciwc}, \campo{crwc}, \campo{cswc} e gli integrali di colonna per lo stesso denominatore,
$1 + r_t$, dove $r_t$ è il rapporto di mescolanza dell'acqua totale. Sul campo tridimensionale la
correzione vale lo \SI{0.011}{\percent} in mediana, l'\SI{1.44}{\percent} al 99° percentile e il
\SI{2.26}{\percent} al massimo.

\paragraph{Frazione nuvolosa.} Con lo schema di Xu--Randall la frazione nuvolosa \campo{cc} è
continua. Su un sottodominio di $500 \times 600$ celle solo fra il \SI{18}{\percent} e il
\SI{22}{\percent} dei punti nuvolosi vale esattamente 1, e gli altri coprono tutto l'intervallo fra 0 e 1.

\paragraph{Umidità relativa.} Il campo \campo{r} è l'umidità relativa di ECMWF, calcolata sulla fase
mista con le equazioni 7.5, 7.6, 7.89 e 7.90 della documentazione IFS Cy41r2, il ciclo con cui è
stato prodotto ERA5 \parencite{era5doc}, e con le costanti di quel ciclo \parencite{ifs41r2p4}. La pressione di vapore
saturo è
\begin{equation}
  e_\text{sat}(T) = \alpha(T)\, e_\text{sat,w}(T) + \bigl(1 - \alpha(T)\bigr)\, e_\text{sat,i}(T),
  \qquad
  \alpha(T) = \left(\frac{T - T_\text{ice}}{T_0 - T_\text{ice}}\right)^{2},
  \label{eq:fase-mista}
\end{equation}
dove $T$ è la temperatura, $e_\text{sat,w}$ ed $e_\text{sat,i}$ sono le pressioni di vapore saturo sull'acqua e sul
ghiaccio (formula di Tetens), $T_\text{ice} = \SI{250.16}{K}$, $T_0 = \SI{273.16}{K}$, e $\alpha$ vale
0 sotto $T_\text{ice}$ e 1 sopra $T_0$. Sotto $T_\text{ice}$ questa umidità relativa differisce fino a
un fattore \num{1.8} da quella calcolata sull'acqua liquida. Il campo non è troncato a \SI{100}{\percent}:
la sovrasaturazione è uno stato del modello, e anche ERA5 la pubblica. Nel file orario del 15 luglio 2024 alle 12 UTC
\campo{r} raggiunge il \SI{139.7}{\percent} a \SI{300}{\hPa} nel regime del ghiaccio, e \num{34785} punti
superano il \SI{100}{\percent} a \SI{400}{\hPa}. Un sistema di acquisizione che scarta valori
sopra il \SI{100}{\percent} scarta dati validi.

\paragraph{Geopotenziale.} Il geopotenziale pubblicato è $z = g_0\, z_\text{WRF}$, con
$g_0 = \SI{9.80665}{m/s^2}$ la gravità standard e $z_\text{WRF}$ l'altezza geopotenziale del
modello. Dividendo \campo{z} per $g_0$ si ritrova esattamente l'altezza del modello, in superficie
e su ogni livello di pressione, e il geopotenziale di superficie e quello dei livelli seguono la
stessa convenzione, verificata punto per punto sulle \num{1349910} celle di terra con quota sopra
\SI{1}{m} del 15 luglio 2024 alle 12 UTC. WRF però integra l'equazione idrostatica con
$g = \SI{9.81}{m/s^2}$, e l'archivio si discosta quindi dello \SI{0.034}{\percent} dall'equilibrio
idrostatico rispetto alle proprie \campo{t} e pressione: una parte su \num{2900}, circa \SI{2}{m} a
\SI{500}{\hPa} e \SI{7}{m} a \SI{50}{\hPa} \parencite{issue39}. Le due proprietà, \campo{z} diviso per $g_0$
uguale all'altezza del modello ed equilibrio idrostatico con \campo{t} e pressione del modello, non
si possono avere insieme.

\section{Sotto il terreno}
\label{sec:sotto-terreno}

Dove un livello di pressione sta sotto il terreno, l'archivio pubblica il valore del livello del
modello più basso, così com'è, senza estrapolare né riempire (\emph{clamping}). Sotto il terreno
\campo{t} coincide esattamente con il livello del modello più basso, \campo{q} entro
\SI{1e-9}{kg/kg} e \campo{z} entro \SI{2.5e-4}{m} \parencite{issue44}. Il valore è uno stato reale
del modello, preso fra \SI{24}{m} e \SI{27}{m} sopra il terreno, e resta coerente con i campi vicino alla superficie. A luglio
$\campo{2t} - \campo{t}(\SI{1000}{\hPa})$ sotto il terreno vale \SI{-0.9}{K} a mezzanotte e
\SI{+1.4}{K} a mezzogiorno, cioè un'inversione notturna e uno strato superadiabatico diurno nei circa
\SI{22}{m} fra il livello di \SI{2}{m} e il livello del modello più basso. Sotto il terreno
\campo{r} arriva a circa il \SI{163}{\percent}, perché accoppia l'umidità del livello del modello
più basso con la pressione nominale di \SI{1000}{\hPa}.

Il prezzo è che \campo{z} contraddice \campo{sp}. Dove la pressione di superficie dice che
\SI{1000}{\hPa} è sotto il terreno, il geopotenziale pubblicato a \SI{1000}{\hPa} lo mette in mediana
fra \SI{24}{m} e \SI{27}{m} sopra. Succede sul \SI{47.6}{\percent} del dominio a \SI{1000}{\hPa} e sul
\SI{9}{\percent} a \SI{925}{\hPa}. ERA5 invece estrapola sotto il terreno con un gradiente termico,
e i due archivi differiscono per convenzione proprio dove \SI{1000}{\hPa} è sotto il terreno. Da qui
viene tutta la differenza a \SI{1000}{\hPa}: sul campione di verifica il bias di \campo{t} rispetto a
ERA5 è \SI{-1.71}{K} a \SI{1000}{\hPa} contro \SI{-0.05}{K} a \SI{850}{\hPa}, e la correlazione di
\campo{z} con ERA5 a \SI{1000}{\hPa} è \num{0.108}. Per un confronto con ERA5 o per calcolare spessori,
i livelli sotto il terreno vanno mascherati con \campo{sp}, che è nello stesso file orario.

\section{Il livello di \SI{50}{\hPa}}
\label{sec:livello-50}

Il top del modello è a \SI{50}{\hPa} (\cref{sec:dominio}), e il livello del modello più alto dove
WRF definisce la temperatura sta sempre sotto, per esempio a \SI{51.9}{\hPa} in media sul dominio il
19 gennaio 2024 alle 00 UTC. Il livello di pressione di \SI{50}{\hPa} sta quindi sopra il livello
del modello più alto, e l'interpolazione
verticale vi ripete i valori di quel livello, come fa il clamping sotto il terreno
(\cref{sec:sotto-terreno}). In quell'istante \campo{t} a \SI{50}{\hPa} coincide con la temperatura
del livello del modello più alto entro la precisione di codifica, \SI{0.009}{K}, e $\campo{z}/g_0$
con la sua altezza geopotenziale entro \SI{1}{m}; lo stesso vale per gli altri campi a
\SI{50}{\hPa}. Nello stesso istante lo strato fra il livello del modello più alto e
\SI{50}{\hPa} ha uno spessore idrostatico di \SI{227}{m} in media, fra \SI{214}{m} e \SI{239}{m} sul dominio, e di tanto
$\campo{z}/g_0$ a \SI{50}{\hPa} sta sotto la quota del livello di pressione di \SI{50}{\hPa}. Lo
spessore fra \SI{100}{\hPa} e \SI{50}{\hPa} calcolato dall'archivio, \SI{4009}{m} in media, è quindi
più corto di circa il \SI{5.4}{\percent} dello spessore idrostatico, circa \SI{4236}{m}. Un vento
geostrofico calcolato con \campo{z} a \SI{50}{\hPa} risente della variazione orizzontale dello
spessore mancante, \SI{25}{m} fra il minimo e il massimo sul dominio.

\section{Campi vicino alla superficie}
\label{sec:vicino-superficie}

La \cref{tab:var-quota} elenca i campi a quota fissa sopra il suolo e la pressione al livello del
mare.

\begin{table}[htbp]
  \centering
  \small
  \caption{Campi a quota fissa e pressione al livello del mare.}
  \label{tab:var-quota}
  \begin{tabular}{@{}lrlll@{}}
    \toprule
    Campo & paramId & Descrizione & Unità & Quota \\
    \midrule
    \campo{10u}, \campo{10v}   & 165, 166       & componenti del vento & \si{m/s} & \SI{10}{m} \\
    \campo{100u}, \campo{100v} & 228246, 228247 & componenti del vento & \si{m/s} & \SI{100}{m} \\
    \campo{200u}, \campo{200v} & 228239, 228240 & componenti del vento & \si{m/s} & \SI{200}{m} \\
    \campo{10fg} & 49     & massimo orario del vento risolto & \si{m/s} & \SI{10}{m} \\
    \campo{2t}   & 167    & temperatura & \si{K} & \SI{2}{m} \\
    \campo{2d}   & 168    & temperatura di rugiada & \si{K} & \SI{2}{m} \\
    \campo{2r}   & 260242 & umidità relativa & \si{\percent} & \SI{2}{m} \\
    \campo{vwsh} & 260068 & shear verticale del vento fra 10 e \SI{100}{m} & \si{s^{-1}} & \SI{100}{m} \\
    \campo{msl}  & 151    & pressione al livello del mare & \si{Pa} & livello del mare \\
    \bottomrule
  \end{tabular}
\end{table}

\paragraph{Umidità relativa a \SI{2}{m}.} \campo{2r} è l'umidità relativa sull'acqua liquida a ogni
temperatura, la convenzione dell'Organizzazione Meteorologica Mondiale per le misure di umidità
in capannina \parencite{wmo8}, cioè per le osservazioni con cui si confronta un'umidità a
\SI{2}{m}. ERA5 non pubblica
un'umidità relativa a \SI{2}{m}. \campo{r} e \campo{2r} seguono quindi due definizioni diverse e
differiscono dove la temperatura è sotto $T_0$. Un profilo che mette \campo{2r} alla base di \campo{r}, o un gradiente
verticale di umidità calcolato fra i due, mescola due grandezze. Per averle nella stessa
convenzione si ricalcola l'umidità relativa da \campo{q}, \campo{t} e \campo{sp}, tutti pubblicati.

\paragraph{Temperatura di rugiada.} L'inversione che ricava la temperatura di rugiada dal vapore dà,
nei punti saturi, un \campo{2d} che supera \campo{2t} di \SI{0.0133}{K}, lo stesso in ogni istante e
in ogni stagione. In quattro istanti di gennaio 2024, luglio 2024 e
febbraio 2025 i punti sono fra \num{6161} e \num{64005}, tutti con umidità relativa del
\SI{100}{\percent} \parencite{issue42}. Il valore è lasciato così. Troncare \campo{2d} a \campo{2t} toglierebbe gli incroci
visibili ma lascerebbe lo stesso scarto in ogni punto non saturo, introdurrebbe un gomito nel campo
al \SI{100}{\percent} e romperebbe la coerenza fra \campo{2d} e il vapore a \SI{2}{m} del modello (\campo{Q2}), che è
entro un errore relativo fra \num{3.6e-4} e \num{1.1e-2} sulla pressione di vapore.
Un'umidità relativa ricostruita da \campo{2t} e \campo{2d} arriva quindi al \SI{100.09}{\percent}.

\section{Colonne e nubi}
\label{sec:colonne}

La \cref{tab:var-colonna} elenca gli integrali di colonna e le coperture nuvolose.

\begin{table}[htbp]
  \centering
  \small
  \caption{Integrali di colonna e coperture nuvolose.}
  \label{tab:var-colonna}
  \begin{tabular}{@{}lrll@{}}
    \toprule
    Campo & paramId & Descrizione & Unità \\
    \midrule
    \campo{tcwv} & 137    & vapore d'acqua totale in colonna & \si{kg/m^2} \\
    \campo{tclw} & 78     & acqua liquida di nube in colonna & \si{kg/m^2} \\
    \campo{tciw} & 79     & ghiaccio di nube in colonna & \si{kg/m^2} \\
    \campo{tcrw} & 228089 & pioggia in colonna & \si{kg/m^2} \\
    \campo{tcsw} & 228090 & neve in colonna & \si{kg/m^2} \\
    \campo{tcw}  & 136    & acqua totale in colonna & \si{kg/m^2} \\
    \campo{tcc}  & 164    & copertura nuvolosa totale & 0--1 \\
    \campo{hcc}  & 188    & copertura nuvolosa alta & 0--1 \\
    \campo{mcc}  & 187    & copertura nuvolosa media & 0--1 \\
    \campo{lcc}  & 186    & copertura nuvolosa bassa & 0--1 \\
    \bottomrule
  \end{tabular}
\end{table}

\paragraph{Acqua totale.} ECMWF definisce \campo{tcw} come la somma di cinque specie, vapore,
acqua liquida, ghiaccio, pioggia e neve \parencite{paramdb}. L'archivio somma sei specie, perché include il graupel
\parencite{issue47}. ECMWF definisce un contenuto di graupel (paramId 260028, \campo{grle}) e una
sua colonna (260001, \campo{tcolg}) \parencite{paramdb}, ma ERA5 non pubblica nessuno dei due, e il graupel sta solo
dentro \campo{tcw}. Ricostruire \campo{tcw} come $\campo{tcwv} + \campo{tclw} + \campo{tciw} +
\campo{tcrw} + \campo{tcsw}$ non lo riproduce, e il residuo è la colonna di graupel. Nei quattro
istanti esaminati, il 15 gennaio, aprile, luglio e ottobre 2024 alle 12 UTC, il residuo è in media
di dominio lo \SI{0.4}{\percent} del condensato. Il suo massimo è \SI{5.69}{kg/m^2} in gennaio,
\SI{13.89}{kg/m^2} in aprile, \SI{16.73}{kg/m^2} in ottobre e \SI{26.46}{kg/m^2} in luglio, in una
cella convettiva in cui supera tutto il resto del condensato della colonna.

\paragraph{Copertura nuvolosa.} \campo{tcc}, \campo{hcc}, \campo{mcc} e \campo{lcc} sono
sovrapposizioni maximum-random della frazione nuvolosa \campo{CLDFRA} di WRF. Sul campione di
verifica, con i campi CHAPTER aggregati sulla griglia ERA5, il rapporto delle medie è \num{0.783}
per la copertura totale e \num{0.610} per quella alta (\cref{tab:nubi}), e la copertura totale
CHAPTER è più bassa di quella ERA5 in ognuno dei 34 istanti \parencite{issue48}.

\begin{table}[htbp]
  \centering
  \small
  \caption{Coperture nuvolose e ghiaccio di nube in colonna (\si{kg/m^2}): medie sul campione di
  verifica, sulla griglia ERA5.}
  \label{tab:nubi}
  \begin{tabular}{@{}lSSS@{}}
    \toprule
    Campo & {CHAPTER} & {ERA5} & {Rapporto delle medie} \\
    \midrule
    \campo{tcc}  & 0.381  & 0.487  & 0.783 \\
    \campo{hcc}  & 0.169  & 0.276  & 0.610 \\
    \campo{mcc}  & 0.161  & 0.205  & 0.786 \\
    \campo{lcc}  & 0.231  & 0.268  & 0.862 \\
    \campo{tciw} & 0.0185 & 0.0166 & 1.109 \\
    \bottomrule
  \end{tabular}
\end{table}

Il divario non viene dall'ipotesi di sovrapposizione. Calcolata punto per punto sulla
\campo{CLDFRA} dell'archivio, l'intera gamma delle ipotesi possibili, dalla sovrapposizione massima
a quella casuale, è larga fra il \SI{2}{\percent} e il \SI{4}{\percent}, mentre la copertura totale
ERA5 supera quella CHAPTER del \SI{28}{\percent}: l'ipotesi spiega al più un settimo del divario.
Il divario sta nella \campo{CLDFRA} stessa, il diagnostico di Xu--Randall. Il ghiaccio di nube in
colonna vale \num{1.109} volte quello di ERA5 mentre \campo{hcc} ne vale \num{0.610}, e lo schema
attribuisce quindi al condensato presente una frazione nuvolosa minore di quella di ERA5. Una parte
del divario viene dalla risoluzione. La copertura CHAPTER è sovrapposta a \SI{3}{km} e poi
aggregata, quella di ERA5 è sovrapposta su un profilo già liscio a \SI{31}{km}, e a parità di
frazioni per strato dà una copertura maggiore. Anche questo effetto è limitato dalla stessa gamma,
fra il \SI{2}{\percent} e il \SI{4}{\percent}.

\section{Instabilità}
\label{sec:instabilita}

\campo{mucape} (paramId 228235) e \campo{mucin} (228236), in \si{J/kg}, sono l'energia potenziale
convettiva disponibile (CAPE) e l'inibizione convettiva della particella più instabile, cioè della
particella mescolata su \SI{500}{m} con la massima temperatura potenziale equivalente nei
\SI{3}{km} più bassi. Non sono calcolate dalla superficie.

La routine che calcola le due grandezze restituisce un'inibizione solo nelle colonne con CAPE di
almeno \SI{100}{J/kg}, e altrove \campo{mucin} è un punto mancante (bitmap). Uno zero in quelle
colonne direbbe che nulla frena la convezione proprio dove la frena di più, perché una colonna
senza CAPE è normalmente una colonna con molta inibizione: dove \campo{mucin} è mancante, ERA5 porta
un'inibizione media fra \SI{25}{J/kg} e \SI{265}{J/kg}. Riempire i punti mancanti con zeri
insegnerebbe a un modello statistico che le situazioni con poca CAPE non hanno inibizione, cioè il
contrario della fisica.

La frazione di dominio in cui \campo{mucin} è presente va dal \SI{3.8}{\percent} in gennaio al
\SI{28.8}{\percent} in luglio (\cref{tab:mucin}). Dove è presente, \campo{mucin} vale fra \num{0.03}
e \num{0.24} volte quello di ERA5, perché le colonne in cui esiste sono le meno inibite.

\begin{table}[htbp]
  \centering
  \small
  \caption{Frazione del dominio in cui \campo{mucin} è presente, alle 12 UTC.}
  \label{tab:mucin}
  \begin{tabular}{@{}lSS@{}}
    \toprule
    Istante di validità & {Presente (\si{\percent})} & {Mancante (\si{\percent})} \\
    \midrule
    2024-01-15 & 3.8  & 96.2 \\
    2024-02-15 & 5.7  & 94.3 \\
    2024-03-15 & 4.8  & 95.2 \\
    2024-07-15 & 28.8 & 71.2 \\
    2024-10-15 & 15.2 & 84.8 \\
    2025-02-15 & 5.5  & 94.5 \\
    \bottomrule
  \end{tabular}
\end{table}

\campo{mucape} è invece denso, senza punti mancanti: dove la routine non trova CAPE, lo zero
pubblicato è il valore corretto. Vale zero su \num{1536450} dei \num{2220273} punti
(\SI{69.2}{\percent}) il 15 febbraio 2024 e su \num{1020809} (\SI{46.0}{\percent}) il 15 luglio
2024, entrambi alle 12 UTC, e su \num{744745} (\SI{33.5}{\percent}) il 4 settembre 2024 alle 15 UTC.

\chapter{Radiazione, temperatura di pelle e scambi di superficie}
\label{cap:superficie}

Sul campione di verifica l'archivio porta più radiazione solare discendente in superficie di ERA5
(\num{1.136} volte), meno copertura nuvolosa totale (\num{0.783} volte), e una temperatura di pelle
sulla terra più fredda di \SI{1.64}{K}. Radiazione, nubi e bilancio energetico della superficie vanno
quindi letti insieme.

\section{Radiazione}
\label{sec:radiazione}

I 13 flussi radiativi della \cref{tab:radiazione} sono campi accumulati sull'ora che termina
all'istante di validità, in \si{J/m^2} (\cref{sec:tempo}).

\begin{table}[htbp]
  \centering
  \small
  \caption{Flussi radiativi, accumulati sull'ora che termina all'istante di validità.}
  \label{tab:radiazione}
  \begin{tabular}{@{}lrl@{}}
    \toprule
    Campo & paramId & Descrizione \\
    \midrule
    \campo{ssrd}  & 169    & radiazione solare discendente in superficie \\
    \campo{ssrdc} & 228129 & radiazione solare discendente in superficie, cielo sereno \\
    \campo{ssr}   & 176    & radiazione solare netta in superficie \\
    \campo{ssrc}  & 210    & radiazione solare netta in superficie, cielo sereno \\
    \campo{strd}  & 175    & radiazione termica discendente in superficie \\
    \campo{strdc} & 228130 & radiazione termica discendente in superficie, cielo sereno \\
    \campo{str}   & 177    & radiazione termica netta in superficie \\
    \campo{strc}  & 211    & radiazione termica netta in superficie, cielo sereno \\
    \campo{tisr}  & 212    & radiazione solare incidente al top dell'atmosfera \\
    \campo{tsr}   & 178    & radiazione solare netta al top dell'atmosfera \\
    \campo{tsrc}  & 208    & radiazione solare netta al top dell'atmosfera, cielo sereno \\
    \campo{ttr}   & 179    & radiazione termica netta al top dell'atmosfera \\
    \campo{ttrc}  & 209    & radiazione termica netta al top dell'atmosfera, cielo sereno \\
    \bottomrule
  \end{tabular}
\end{table}

La radiazione solare in superficie è più alta di quella di ERA5: la discendente
dell'\SI{8.5}{\percent} in cielo sereno e del \SI{13.6}{\percent} con le nubi (\cref{tab:rad-era5})
\parencite{issue45}. La simulazione non ha aerosol (\campo{aer_opt=0}), mentre ERA5 usa la
climatologia degli aerosol di \textcite{tegen1997}. Il flusso in cielo sereno isola l'atmosfera
senza nubi, e i campi termici, su cui l'aerosol non agisce, concordano con ERA5 entro
l'\SI{1.1}{\percent}: il confronto attribuisce all'assenza di aerosol l'eccesso in cielo sereno.
Con le nubi l'eccesso cresce, perché la copertura nuvolosa totale vale \num{0.783} volte quella di
ERA5 (\cref{sec:colonne}) e lascia passare più radiazione. Per applicazioni di energia solare o per l'addestramento contro
osservazioni di piranometro l'archivio porta un eccesso sistematico di questa entità, che è una
proprietà della simulazione e che la conversione non può togliere.

\begin{table}[htbp]
  \centering
  \small
  \caption{Rapporto delle medie fra CHAPTER ed ERA5 sul campione di verifica, con i campi CHAPTER
  aggregati sulla griglia ERA5.}
  \label{tab:rad-era5}
  \begin{tabular}{@{}lS@{}}
    \toprule
    Campo & {CHAPTER / ERA5} \\
    \midrule
    \campo{ssrdc} & 1.085 \\
    \campo{ssrd}  & 1.136 \\
    \campo{ssrc}  & 1.116 \\
    \campo{ssr}   & 1.178 \\
    \campo{tsrc}  & 1.031 \\
    \campo{strd}  & 0.989 \\
    \campo{strdc} & 0.996 \\
    \campo{ttr}   & 1.009 \\
    \campo{ttrc}  & 0.990 \\
    \bottomrule
  \end{tabular}
\end{table}

\section{Temperatura di pelle}
\label{sec:pelle}

La simulazione non ha scritto la temperatura di pelle del modello (\campo{TSK}), e l'archivio la
ricostruisce come \campo{skt} dalla radiazione termica uscente in superficie. Il flusso ascendente
$F^{\uparrow}$ (\campo{LWUPB} in WRF) è
\begin{equation}
  F^{\uparrow} = \varepsilon \sigma T_s^4 + (1 - \varepsilon) F^{\downarrow},
  \label{eq:lwup}
\end{equation}
con $\varepsilon$ l'emissività della superficie, $\sigma$ la costante di Stefan--Boltzmann, $T_s$ la
temperatura di pelle e $F^{\downarrow}$ il flusso termico discendente in superficie (\campo{GLW}). L'emissività
si ricava dalla simulazione stessa. RRTMG calcola anche i flussi
in cielo sereno, con la stessa emissività e la stessa temperatura di pelle e un diverso flusso
discendente. Sottraendo le due forme della \cref{eq:lwup} la temperatura scompare:
\begin{equation}
  \varepsilon = 1 - \frac{F^{\uparrow} - F^{\uparrow}_\text{cs}}{F^{\downarrow} - F^{\downarrow}_\text{cs}},
  \label{eq:emissivita}
\end{equation}
dove il pedice cs indica il cielo sereno. L'espressione è un'identità, senza temperatura di
riferimento e senza parametri stimati. Misurata così, l'emissività di categoria è la stessa in ogni
cella con la stessa categoria dominante, con uno scarto interquartile nullo alla quarta cifra
decimale, e vale \num{0.98000} sull'acqua. I valori per categoria
coincidono con la colonna \campo{LEMI} della sezione \campo{MODI-RUC} di \campo{VEGPARM.TBL}
\parencite{wrf411}, da \num{0.850} sul suolo nudo a \num{0.980} sull'acqua.

Con la neve l'emissività passa a \num{0.98} per tutte le categorie. L'archivio applica \num{0.98}
dove la copertura nevosa \campo{SNOWC} supera la soglia di copertura nevosa, \num{0.015}
(\SI{1.5}{\percent}), e sotto la soglia raccorda i due valori. La soglia di copertura nevosa è
l'unico parametro della ricostruzione non letto dalla simulazione: è stata scelta come il minimo dell'errore sulla
temperatura di pelle in una scansione sui 34 istanti del campione di verifica.

L'accuratezza dipende dalla stagione (\cref{tab:skt}). Sull'acqua aperta, dove la temperatura di
pelle del modello è la temperatura del mare, la ricostruzione la ritrova con un errore quadratico
medio di \SI{0.0005}{K} su circa \num{910000} punti, in ogni istante del campione di verifica. Sulla terra l'emissività
ricostruita coincide con quella del modello entro $10^{-4}$ sul \SI{99.98}{\percent} dei punti a
luglio e sul \SI{93.05}{\percent} a gennaio. Tutta la differenza fra le stagioni viene dalla regola
della neve. Sulla terra senza neve la temperatura di pelle ricostruita differisce dalla temperatura del nodo
del suolo a \SI{0}{cm} (\campo{TSLB}) di \SI{0.026}{K} in errore quadratico medio, il 4 settembre
2024 alle 16 UTC.

\begin{table}[htbp]
  \centering
  \small
  \caption{Accuratezza della temperatura di pelle ricostruita rispetto allo stato del modello.}
  \label{tab:skt}
  \begin{tabularx}{\textwidth}{@{}Xl@{}}
    \toprule
    Grandezza & Valore \\
    \midrule
    acqua aperta, errore quadratico medio (errore assoluto massimo) & \SI{0.0005}{K} (\SI{0.002}{K}) \\
    terra, punti con emissività entro $10^{-4}$ & \SI{99.98}{\percent} a luglio, \SI{93.05}{\percent} a gennaio \\
    terra, errore quadratico medio a luglio & \SI{0.0012}{K} \\
    terra, errore quadratico medio a gennaio e febbraio & da \SI{0.084}{K} a \SI{0.119}{K} \\
    terra, gennaio e febbraio, errore assoluto al 99° percentile (massimo) & \SI{0.81}{K} (\SI{2.79}{K}) \\
    \bottomrule
  \end{tabularx}
\end{table}

Rispetto a ERA5 la temperatura di pelle sulla terra è più fredda di \SI{1.64}{K} in media sui 34
istanti del campione di verifica. Lo scarto, CHAPTER meno ERA5, va da \SI{-0.5}{K} a \SI{-1.2}{K}
fra gli istanti delle 00 UTC\todo{istanti delle 00 UTC da rimisurare sui file della run del giorno prima (issue 72)} e da \SI{-0.8}{K} a \SI{-3.24}{K} fra quelli delle 12 UTC, con il
valore più freddo il 15 settembre 2024. Sul mare lo scarto è \SI{+0.01}{K}, e la correlazione con
ERA5 è fra \num{0.95} e \num{0.98} sia sulla terra sia sul mare \parencite{issue57}. Sul mare, dove
le due rianalisi usano quasi la stessa temperatura del mare, lo scarto di \SI{0.01}{K} esclude come
causa la ricostruzione, l'emissività e la codifica. L'errore della ricostruzione sulla terra, al più
\SI{0.119}{K} in errore quadratico medio, è più di dieci volte più piccolo dello scarto medio
rispetto a ERA5.

Lo scarto sulla terra va nel verso opposto a quello che suggeriscono radiazione e nubi. Più
radiazione solare al suolo e meno nubi portano energia alla superficie, e la superficie è più fredda,
soprattutto alle 12 UTC, quando l'eccesso di radiazione solare è massimo. La ripartizione
dell'energia in superficie resta quindi non spiegata. Per chiarirla servirebbe
il bilancio energetico di superficie, e i tre campi che lo chiudono, i flussi di calore sensibile e
latente e il flusso di calore nel suolo, non sono nell'archivio (\cref{cap:limiti}). Senza questi
campi non si separano quattro cause possibili: lo schema di superficie può mandare più energia in
calore latente, il clima dell'umidità del suolo può essere diverso da quello di ERA5, l'orografia
risolta a \SI{3}{km} alza la quota media di una cella aggregata sulla griglia ERA5, oppure lo scarto è
un errore del modello. Chi confronta \campo{skt} con temperature di superficie da satellite, o ne
corregge lo scarto, deve usare lo scarto misurato e non una stima ricavata da radiazione e nubi.

\section{Rugosità, velocità d'attrito e sforzo}
\label{sec:rugosita}

La rugosità \campo{fsr} è la rugosità locale \campo{ZNT} di WRF: un valore di tabella per categoria
dominante, sostituito sotto la neve. Vale \SI{0.70}{m} per le foreste, \SI{0.15}{m} per
savane e prati, \SI{0.065}{m} per il suolo nudo, \SI{1.00}{m} sui \num{19867} punti urbani, e scende
a \SI{0.011}{m} dove la neve copre la superficie. La rugosità di ERA5 è invece una rugosità efficace,
che contiene la resistenza orografica che l'IFS parametrizza a \SI{31}{km}; a \SI{3}{km} l'orografia
è risolta. Sul mare le due rugosità seguono entrambe la formula di Charnock. Sulla terra differiscono di un
fattore da due a tre (\cref{tab:rugosita}), e un coefficiente di resistenza calcolato dalla
\campo{fsr} dell'archivio differisce da uno calcolato da quella di ERA5 per la sola definizione della
rugosità \parencite{issue56}.

\begin{table}[htbp]
  \centering
  \small
  \caption{Rugosità e velocità d'attrito: intervallo delle medie per istante su sei istanti
  distribuiti sulle stagioni, e rapporti con ERA5.}
  \label{tab:rugosita}
  \begin{tabularx}{\textwidth}{@{}lXXXX@{}}
    \toprule
    Campo & CHAPTER & ERA5 & Rapporto delle medie & Rapporto delle mediane \\
    \midrule
    \campo{fsr}, terra  & \SIrange{0.188}{0.217}{m}   & \SIrange{0.375}{0.394}{m}   & \numrange{0.48}{0.55} & \numrange{0.32}{0.48} \\
    \campo{fsr}, mare   & \SIrange{0.0089}{0.0091}{m} & \SIrange{0.0103}{0.0105}{m} & \numrange{0.85}{0.88} & \numrange{1.5}{1.9} \\
    \campo{zust}, terra & \SIrange{0.44}{0.52}{m/s}   & \SIrange{0.23}{0.28}{m/s}   & \numrange{1.62}{2.06} & \\
    \campo{zust}, mare  & \SIrange{0.21}{0.34}{m/s}   & \SIrange{0.19}{0.31}{m/s}   & \numrange{1.08}{1.12} & \\
    \bottomrule
  \end{tabularx}
\end{table}

La velocità d'attrito \campo{zust} è quella del modello. Il rapporto delle medie con ERA5 è fra
\num{1.62} e \num{2.06} sulla terra e fra \num{1.08} e \num{1.12} sul mare, dove la correlazione è
fra \num{0.85} e \num{0.90}. Una rugosità fra \num{0.32} e \num{0.55} volte quella di ERA5, insieme a
una velocità d'attrito fra \num{1.62} e \num{2.06} volte, esclude un semplice errore di scala sulla
stessa rugosità. A \SI{3}{km} il vento forzato dal terreno è risolto, e la velocità d'attrito del
modello lo segue.

Le componenti dello sforzo in superficie \campo{iews} e \campo{inss}, in \si{N/m^2}, sono il modulo
$\rho u_*^2$, con $\rho$ la densità dell'aria e $u_*$ la velocità d'attrito del modello, proiettato
sulla direzione del vento a \SI{10}{m}. L'allineamento con il vento a \SI{10}{m} è l'unica ipotesi.
Con $\tau$ il vettore di componenti \campo{iews} e \campo{inss}, le identità $|\tau| = \rho u_*^2$ e
$u_* = \sqrt{|\tau|/\rho}$ valgono entro \SI{1.8e-6}{N/m^2} ed entro \SI{2.7e-4}{m/s} in ogni punto
di sei istanti.

\chapter{Vento a 10 m e raffica}
\label{cap:raffica}

\section{Il campo \texorpdfstring{\campo{10fg}}{10fg}}
\label{sec:10fg}

Il campo \campo{10fg} (paramId 49) è la velocità massima del vento a \SI{10}{m} che la simulazione
ha prodotto nell'ora precedente l'istante di validità. Viene da \campo{WSPD10MAX}, uno dei massimi
orari che WRF scrive con \campo{nwp_diagnostics=1}. È un \emph{vento risolto}: non contiene un
termine turbolento né una stima di ciò che la griglia non rappresenta \parencite{issue50}.

La raffica di ERA5 è un'altra grandezza. L'Organizzazione Meteorologica Mondiale definisce la
raffica come il massimo del vento mediato su \SI{3}{s}, un intervallo più breve di un passo temporale
di un modello a \SI{31}{km}, e ERA5 ne pubblica una stima, la \emph{raffica parametrizzata}
\parencite{paramdb}. Nell'IFS la raffica parametrizzata è il vento medio a \SI{10}{m} più un termine
turbolento proporzionale alla velocità d'attrito, più, in condizioni convettive, un termine
proporzionale allo shear del vento nello strato \SIrange{850}{950}{\hPa} \parencite[sez.~3.10.4,
eq.~3.99]{ifs41r2p4}. L'IFS la calcola a ogni passo temporale e ne archivia il massimo dall'ultima
scrittura.

La simulazione regionale riceve da ERA5 lo stato dell'atmosfera e della superficie, non i campi
diagnostici. La tabella standard dei campi ECMWF di WPS 4.1 (\campo{Vtable.ECMWF}) non contiene né il
paramId 49 né il 228029 \parencite{wps41}; la tabella usata da questa simulazione non è stata
verificata. Fra le 200 variabili di un file wrfout nessun campo di raffica si aggiunge a
\campo{WSPD10MAX}, e il campo pubblicato come \campo{10fg} è il più vicino a una raffica che la
simulazione offra.

\section{Confronto con ERA5}
\label{sec:10fg-era5}

La \cref{tab:10fg} riporta il confronto su 16 istanti di validità. \campo{10fg} vale 0,625 volte il
\campo{10fg} di ERA5 e 0,644 volte la raffica parametrizzata istantanea \campo{i10fg}, con
correlazioni di 0,813 e 0,820. Rispetto al vento istantaneo a \SI{10}{m} dello stesso file vale fra
1,004 e 1,032 volte, e in nessuno dei \num{2220273} punti è più basso. Sotto il paramId 49
l'archivio pubblica la velocità del vento a \SI{10}{m}, mentre la raffica parametrizzata di ERA5 vale
circa 1,6 volte il suo vento medio. Il rapporto con il \campo{10fg} di ERA5 resta fra 0,593 e 0,673 in
tutti i 34 istanti del campione di verifica, in ogni mese e a ogni ora, ed è una
differenza di definizione.

\begin{table}[htbp]
  \centering
  \small
  \caption{\campo{10fg} dell'archivio rispetto a tre riferimenti, su 16 istanti di validità.}
  \label{tab:10fg}
  \begin{tabular}{@{}lll@{}}
    \toprule
    Riferimento & CHAPTER / riferimento & Correlazione \\
    \midrule
    \campo{10fg} di ERA5 (raffica parametrizzata, massimo orario) & 0,625 & 0,813 \\
    \campo{i10fg} di ERA5 (raffica parametrizzata, istantanea) & 0,644 & 0,820 \\
    vento istantaneo a \SI{10}{m} di CHAPTER, stesso file & 1,004--1,032 & \\
    \bottomrule
  \end{tabular}
\end{table}

A \SI{3}{km} una parte di ciò che ERA5 deve parametrizzare è risolta, come mostra la velocità
d'attrito sulla terra, fra 1,62 e 2,06 volte quella di ERA5 (\cref{sec:rugosita}). Letto come
raffica, \campo{10fg} sottostima di circa il \SI{38}{\percent}, e nessuna correzione del bias lo rende
una raffica, perché è una grandezza diversa. Con una correlazione di 0,813 con il \campo{10fg} di
ERA5 il campo si presta a ordinare luoghi e ore per ventosità e a descrivere la struttura spaziale
del vento, e fa da limite inferiore della raffica; non si presta a stimare il superamento di una
soglia di raffica.

\section{La finestra oraria}
\label{sec:10fg-finestra}

I messaggi di \campo{10fg} dichiarano un massimo sull'ora precedente (\cref{sec:tempo}). Anche la
raffica parametrizzata oraria di ERA5 è un massimo sull'ora precedente, e la lunghezza della
finestra non è la differenza fra i due campi. La codifica è corretta. WRF azzera \campo{WSPD10MAX}
alla scrittura di ogni file wrfout, ma su sei giorni interi, il 15 luglio 2019 e il 1º e 15 febbraio,
20 marzo, 15 luglio e 15 agosto 2024, per 138 passaggi da un'ora alla successiva, l'azzeramento è
mancato 17 volte, il \SI{12.3}{\percent} dei casi \parencite{issue49}. In quelle ore il campo
contiene il massimo su due ore o più, sotto un messaggio che ne dichiara una.

L'effetto si vede nell'eccesso medio rispetto al vento istantaneo a \SI{10}{m}:
\SIrange{0.068}{0.311}{m/s} nelle ore azzerate, \SIrange{0.499}{0.784}{m/s} in quelle non
azzerate. Su una media di dominio di circa \SI{4.5}{m/s} è una sovrastima di
\SIrange{10}{17}{\percent}. Nelle ore non azzerate il campo si avvicina alla raffica parametrizzata di
ERA5 perché la finestra è più lunga, e il valore non permette di riconoscerle. Neppure gli altri campi
del file lo permettono: i sette massimi orari di WRF si azzerano con lo stesso evento e, quando
l'azzeramento manca, restano tutti insieme, per cui nessuno di essi segnala il mancato azzeramento
degli altri. Riconoscerle richiede il file wrfout dell'ora precedente, che al momento della
conversione non è disponibile. Il meccanismo del mancato azzeramento non è stato stabilito.

Nelle statistiche su una stagione il \SI{12.3}{\percent} delle ore porta una finestra più lunga di
quella dichiarata; per un'ora singola la finestra dichiarata di \campo{10fg} non è garantita.

\chapter{Acqua, suolo, neve e campi statici}
\label{cap:suolo}

\section{Precipitazione e deflusso}
\label{sec:precipitazione}

La precipitazione e il deflusso sono campi accumulati sull'ora che termina all'istante di validità
(\cref{sec:tempo}). La
precipitazione totale \campo{tp} (paramId 228, \si{m}) è la precipitazione risolta \campo{RAINNC} di
WRF, completa perché la precipitazione convettiva \campo{RAINC} è identicamente nulla
(\campo{CU_PHYSICS=0}, \cref{tab:configurazione}). La neve \campo{sf} (144, \si{m} di acqua
equivalente) è la somma della neve e del graupel risolti, $\campo{SNOWNC} + \campo{GRAUPELNC}$, cioè
tutta la precipitazione solida come nella definizione ECMWF \parencite{paramdb}. La pioggia
\campo{tirf} (235015, \si{kg/m^2}) è la precipitazione totale meno neve e graupel; la microfisica non
ha grandine. I tre campi si chiudono: con \campo{tirf} convertito in metri,
$\campo{tp} - \campo{sf} - \campo{tirf}$ vale al più \SI{7.8e-8}{m} nel punto peggiore, il 15
febbraio e il 15 luglio 2024 alle 12 UTC.

I tre deflussi, superficiale (\campo{sro}, paramId 8), sotto-superficiale (\campo{ssro}, 9) e totale
(\campo{ro}, 205), portano l'informazione di un campo solo, e l'archivio li pubblica tutti e tre
\parencite{issue51}. In otto istanti del 2024 (quattro nella \cref{tab:deflusso}), il 15 gennaio,
febbraio, marzo e agosto alle 12 UTC e il 15 luglio alle 00, 06, 12 e 18 UTC, \campo{ssro} è
identicamente nullo in quattro e, dove non lo è, è diverso da zero in al più 17 dei \num{2220273}
punti di griglia. In ogni punto degli otto istanti, $|\campo{ro} - \campo{sro} - \campo{ssro}|$ non
supera \SI{7.5e-9}{m}.

\begin{table}[htbp]
  \centering
  \small
  \caption{Medie di dominio dei deflussi, accumulati sull'ora che termina all'istante di validità, in
  \si{m}, e numero di punti di griglia in cui \campo{ssro} è diverso da zero.}
  \label{tab:deflusso}
  \begin{tabular}{@{}lSSrS@{}}
    \toprule
    Istante di validità & {\campo{sro}} & {\campo{ssro}} & Punti \campo{ssro} $\neq 0$ & {\campo{ro}} \\
    \midrule
    2024-02-15 12 UTC & 4.819e-6 & 0       & 0  & 4.819e-6 \\
    2024-07-15 12 UTC & 5.543e-6 & 2.2e-9  & 3  & 5.545e-6 \\
    2024-07-15 18 UTC & 1.829e-5 & 1.5e-9  & 2  & 1.829e-5 \\
    2024-08-15 12 UTC & 2.334e-6 & 0       & 0  & 2.334e-6 \\
    \bottomrule
  \end{tabular}
\end{table}

\campo{ssro} è l'incremento orario di \campo{UDROFF}, l'accumulatore del deflusso sotto-superficiale
del modello. L'accumulatore è diverso da zero in un numero di punti fra 462 e
\num{8363}, con massimi fra 11 e \SI{106}{mm}. Nella run partita alle 18 UTC del 14 luglio 2024 gli
stessi \num{1647} punti portano deflusso sotto-superficiale alla sesta, dodicesima, diciottesima e
ventiquattresima ora (00, 06, 12 e 18 UTC del 15 luglio), solo 19 di essi crescono fra la sesta e la
ventiquattresima ora, e in una singola ora ne crescono al più 17.

Rispetto a ERA5, con le accumulazioni di entrambe le rianalisi sommate sulle ore dalle 01 UTC all'istante di validità, \campo{sro} vale fra \num{1.0} e \num{4.6} volte quello di ERA5, \campo{ssro} fra
\num{0.000} e \num{0.004} volte e \campo{ro} fra \num{0.19} e \num{0.85} volte. In ERA5 il termine
maggiore è il deflusso sotto-superficiale, che nell'archivio è quasi ovunque nullo
(\cref{tab:deflusso}).

\section{Suolo}
\label{sec:suolo}

L'umidità volumetrica del suolo \campo{swvl1-4} (paramId 39--42, \si{m^3/m^3}) e la temperatura del
suolo \campo{stl1-4} (139, 170, 183, 236, \si{K}) sono date sui quattro strati di ERA5, 0--7, 7--28,
28--100 e \SIrange{100}{289}{cm}. Lo schema RUC porta invece valori puntuali su sei nodi, a
\SIlist{0;5;20;40;160;300}{cm}, con un profilo lineare fra un nodo e il successivo. I quattro strati
ERA5 stanno fra il primo e l'ultimo nodo, perché ERA5 comincia a \SI{0}{cm}, dove sta il primo nodo,
e finisce a \SI{289}{cm}, prima del nodo a \SI{300}{cm}. Ogni strato è quindi la media esatta, sullo
strato, del profilo lineare del modello, senza estrapolazione a nessuna delle due estremità. Ogni
strato è una combinazione pesata di tre nodi consecutivi, a 0, 5 e \SI{20}{cm}, a 5, 20 e
\SI{40}{cm}, a 20, 40 e \SI{160}{cm} e a 40, 160 e \SI{300}{cm}, con pesi che sommano a 1.
\campo{swvl} usa l'umidità totale del suolo, liquida più ghiacciata, che è ciò che il nome ERA5
indica.

Gli otto campi sono definiti sulla terra permanente del file statico, e il ghiaccio marino non ha
colonna di suolo, come in ERA5, la cui maschera terra--mare è permanente. Dove si forma ghiaccio
marino WRF riclassifica la cella come terra, le assegna il tipo di suolo 16 e riempie d'acqua la
colonna di suolo: \campo{swvl} arriva esattamente a \num{1.000}, contro una porosità di \num{0.435}
per quel tipo di suolo. Quelle colonne sono uno stato del modello ma non un suolo, e l'archivio le
maschera. Il ghiaccio stesso resta nel file, in \campo{ci}.

\section{Neve}
\label{sec:neve}

L'archivio porta l'acqua equivalente della neve \campo{sd} (paramId 141, \si{m}), la densità
\campo{rsn} (33, \si{kg/m^3}), la copertura \campo{snowc} (260038, \si{\percent}) e la temperatura
dello strato di neve \campo{tsn} (238, \si{K}). \campo{tsn} è \campo{SOILT1}, la temperatura del
primo livello della colonna di neve e suolo del modello, e dove la neve non copre la cella quel
livello è suolo. Il campo è quindi definito solo dove \campo{SNOWC} supera \num{0.9}, cioè dove
\campo{snowc} supera il \SI{90}{\percent}, e altrove è mancante, comprese le celle con neve su una
frazione minore. Il campo è interamente mancante in 7 dei 34 istanti del campione di verifica, tutti
estivi.

\section{Vegetazione, tipo di suolo e laghi}
\label{sec:vegetazione}

\paragraph{Coperture.} \campo{cvl} (paramId 27) e \campo{cvh} (28) sono le somme delle frazioni del
file statico sulle classi MODIS-IGBP di vegetazione bassa e alta, la stessa grandezza intesa da ERA5
(\cref{tab:tipi-veg}). Sono vegetazione alta le classi IGBP con copertura arborea sopra il
\SI{30}{\percent}, cioè le foreste (sopra il \SI{60}{\percent}) e le savane alberate (fra il 30 e il
\SI{60}{\percent}); la savana, fra il 10 e il \SI{30}{\percent}, è bassa \parencite{modis2022}.

\paragraph{Tipi di vegetazione.} \campo{tvl} (29) e \campo{tvh} (30) sono, fra le classi basse e fra
quelle alte, la classe con la frazione maggiore nella cella, tradotta nella tabella di codici WMO
4.234 con la corrispondenza della \cref{tab:tipi-veg}. La tabella distingue arbusti sempreverdi
(codice 16) e decidui (17), mentre MODIS non registra la fenologia degli arbusti. Sulle due classi di
arbusteto il codice non si può scegliere senza un'ipotesi, e \campo{tvl} è mancante sulle
\num{124672} celle in cui domina un arbusteto, il \SI{16.02}{\percent} delle \num{778381} celle con
$\campo{cvl} > 0$. \campo{cvl} riporta comunque la copertura in quelle celle.

\begin{table}[htbp]
  \centering
  \small
  \caption{Traduzione delle classi MODIS-IGBP nei codici WMO 4.234 per \campo{tvl} e \campo{tvh}.}
  \label{tab:tipi-veg}
  \begin{tabularx}{\textwidth}{@{}XlX@{}}
    \toprule
    Classe MODIS-IGBP & Insieme & Codice 4.234 \\
    \midrule
    1 foresta di aghifoglie sempreverdi & alta & 3 aghifoglie sempreverdi \\
    2 foresta di latifoglie sempreverdi & alta & 6 latifoglie sempreverdi \\
    3 foresta di aghifoglie decidue & alta & 4 aghifoglie decidue \\
    4 foresta di latifoglie decidue & alta & 5 latifoglie decidue \\
    5 foresta mista & alta & 18 foresta mista \\
    8 savana alberata & alta & 19 foresta interrotta \\
    9 savana & bassa & 7 erba alta \\
    10 prateria & bassa & 2 erba bassa \\
    11 zona umida permanente (assente) & bassa & 13 paludi e torbiere \\
    12 coltivi & bassa & 1 coltivi, agricoltura mista \\
    14 mosaico di coltivi e vegetazione naturale (assente) & bassa & 1 coltivi, agricoltura mista \\
    18, 19, 20 tundra alberata, mista e nuda & bassa & 9 tundra \\
    6, 7 arbusteto chiuso e aperto & bassa & nessuno: \campo{tvl} mancante \\
    13, 15, 16, 17, 21 urbano, neve e ghiaccio, suolo nudo, acqua, lago & nessuno & nessuno: frazione $1 - \campo{cvl} - \campo{cvh}$ \\
    \bottomrule
  \end{tabularx}
\end{table}

\paragraph{Tipo di suolo.} \campo{slt} (paramId 43) non traduce categorie. Cinque dei sette tipi di
suolo di ECMWF sono definiti da soglie sulle percentuali di argilla e sabbia \parencite{ifs41r2p4},
e l'archivio le applica alle frazioni di argilla e sabbia del file statico (\cref{tab:slt}). Le
soglie dei tipi grossolano, medio e fine sono quelle della FAO \parencite{hwsd2012}; quelle dei tipi
medio fine e molto fine non hanno una fonte esterna. Dopo le soglie, il tipo 6 sostituisce il
risultato dove la categoria di suolo dominante è il materiale organico. Tutte le \num{1304109} celle
di terra permanente hanno un tipo: \num{1069456} medio, \num{187378} grossolano, \num{45450} fine,
\num{1498} organico e 327 molto fine. I tipi 3 (medio fine) e 7 (organico tropicale) non compaiono.
eccodes decodifica \campo{slt} con la tabella WMO 4.213, che ha 11 classi e un significato diverso; i
valori dell'archivio vanno letti con la definizione ECMWF.

\begin{table}[htbp]
  \centering
  \small
  \caption{Soglie per il tipo di suolo \campo{slt} (argilla e sabbia in percentuale).}
  \label{tab:slt}
  \begin{tabular}{@{}ll@{}}
    \toprule
    Condizione & Tipo ECMWF \\
    \midrule
    argilla $\geq 60$ & 5 molto fine \\
    $35 \leq$ argilla $< 60$ & 4 fine \\
    argilla $< 35$ e sabbia $\leq 15$ & 3 medio fine \\
    $18 \leq$ argilla $< 35$ e sabbia $> 15$, oppure argilla $< 18$ e $15 <$ sabbia $\leq 65$ & 2 medio \\
    argilla $< 18$ e sabbia $> 65$ & 1 grossolano \\
    categoria di suolo dominante di materiale organico & 6 organico \\
    \bottomrule
  \end{tabular}
\end{table}

\paragraph{Laghi.} La copertura dei laghi \campo{cl} (paramId 26) viene dalla frazione della classe
di lago del file statico, ed è diversa da zero su \num{32000} celle. Non è ridondante rispetto alla
maschera terra--mare oraria \campo{lsm}: \num{19309} celle che la maschera conta come terra
contengono un lago più piccolo della cella. La profondità dei laghi \campo{dl} (228007, \si{m}) è il
campo \campo{LAKE_DEPTH} del file statico, e non è uno stato del modello, perché la fisica dei laghi è
spenta (\campo{SF_LAKE_PHYSICS=0}). \campo{dl} è definito sulle stesse \num{32000} celle di
\campo{cl} ed è mancante altrove; in \num{22991} di esse vale \SI{10}{m}, il valore predefinito di
WPS. Nel file statico \campo{LAKE_DEPTH} vale questo valore predefinito sul \SI{99.56}{\percent} del
dominio.

\section{Albedo}
\label{sec:albedo}

L'albedo di fondo \campo{al} (paramId 174) è un valore di tabella per categoria dominante. La
simulazione non usa albedo mensili da satellite, e il valore viene dalle colonne invernale ed estiva
della tabella di uso del suolo: la colonna estiva vale da aprile a ottobre, quella invernale da
novembre a marzo. Su tutto il dominio il campo prende da 8 a 11 valori, e al passaggio fra le due
colonne cambia in \num{1171937} celle, l'\SI{89.9}{\percent} delle celle di terra permanente. Ogni passaggio è una riga
della tabella: per esempio da \num{0.20} a \num{0.17} sui coltivi in \num{497649} celle e da
\num{0.23} a \num{0.25} sul suolo nudo in \num{369219}. Sull'acqua vale \num{0.08}, sul ghiaccio
marino \num{0.65}. Nell'arco di un giorno il valore sulla terra non cambia; cambia solo nelle celle di
ghiaccio marino, che seguono la maschera oraria (\cref{sec:statici-orari}). \campo{al} non è
un'albedo osservata né un'albedo che evolve, e una differenza fra due date a cavallo del cambio di
colonna non è un segnale fisico.

L'albedo effettiva della simulazione è \campo{fal} (paramId 243), il rapporto fra radiazione solare
riflessa e incidente in superficie. È definita solo dove la radiazione incidente supera
\SI{50}{W/m^2}, e altrove è mancante.

\section{Orografia}
\label{sec:orografia}

L'orografia è il geopotenziale di superficie \campo{z} (paramId 129). La pendenza \campo{slor} (163)
è il modulo del gradiente della quota del terreno risolto a \SI{3}{km}, $|\nabla h|$, mentre il
paramId 163 indica la pendenza dell'orografia sottogriglia. eccodes mostra il nome del parametro,
\emph{Slope of sub-gridscale orography}, e nulla nel messaggio dice che si tratta della pendenza
risolta. La deviazione standard dell'orografia \campo{sdor} (160), calcolata da WPS, è la grandezza
di ERA5, valutata su una cella di \SI{3}{km} anziché di \SI{31}{km}. Il rapporto delle mediane
rispetto a ERA5 è \num{0.244}, pari a $(3/31)^{0{,}605}$, e la differenza viene dalla risoluzione.

\chapter{Limiti della corrispondenza}
\label{cap:limiti}

Le grandezze ERA5 della \cref{tab:limiti} non hanno un corrispettivo nell'archivio. Il limite
sta nei file wrfout, che non contengono un campo calcolato dal modello, nella configurazione della
simulazione (\cref{sec:configurazione}) o nella regola di pubblicazione (\cref{sec:pubblicazione}).

\begin{table}[htbp]
  \centering
  \small
  \caption{Grandezze ERA5 assenti dall'archivio e origine del limite.}
  \label{tab:limiti}
  \begin{tabularx}{\textwidth}{@{}XX@{}}
    \toprule
    Parametro & Origine del limite \\
    \midrule
    altezza dello strato limite (\campo{blh}) & file wrfout: \campo{PBLH} calcolato e non scritto \\
    flussi di calore sensibile e latente & file wrfout: \campo{HFX} e \campo{LH} calcolati e non scritti \\
    energia cinetica turbolenta & configurazione: nessuno schema con energia cinetica turbolenta \\
    temperatura massima e minima a \SI{2}{m} (\campo{mx2t}, \campo{mn2t}) & configurazione: \campo{output_diagnostics=0} \\
    onde, coefficiente di resistenza che dipende dallo stato del mare & configurazione: nessun modello d'onda \\
    correnti, altezza della superficie del mare & configurazione: oceano non simulato \\
    spessore e albedo del ghiaccio marino & configurazione: ghiaccio solo come dato di ingresso \\
    precipitazione convettiva (\campo{cp}, \campo{csf}) & configurazione: convezione risolta, campo identicamente nullo \\
    componenti della raffica (\campo{10efg}, \campo{10nfg}) & regola di pubblicazione \\
    indice di area fogliare per classe (\campo{lai_lv}, \campo{lai_hv}) & regola di pubblicazione \\
    angolo e anisotropia dell'orografia sottogriglia (\campo{anor}, \campo{isor}) & regola di pubblicazione \\
    fusione della neve (\campo{smlt}) & regola di pubblicazione \\
    gli 8 livelli di pressione ERA5 sopra \SI{50}{\hPa} (\SIlist{1;2;3;5;7;10;20;30}{\hPa}) & configurazione: top del modello a \SI{50}{\hPa} \\
    \bottomrule
  \end{tabularx}
\end{table}

\section{Campi calcolati e non scritti}
\label{sec:non-scritti}

Lo schema YSU calcola un'altezza dello strato limite (\campo{PBLH}), e con \campo{isfflx=1} i flussi
di calore sensibile e latente (\campo{HFX}, \campo{LH}) sono calcolati a ogni passo temporale. Nessuno
dei tre è stato scritto nei file wrfout, e non esiste un output della simulazione più completo. Lo
stesso vale per la temperatura di pelle del modello \campo{TSK}, che l'archivio ricostruisce dalla
radiazione (\cref{sec:pelle}).

Per l'altezza dello strato limite una diagnosi a posteriori, per esempio con il numero di Richardson
di bulk sui profili, sarebbe nostra e non del modello, e non soddisfa la regola di pubblicazione
(\cref{sec:pubblicazione}). Nemmeno i flussi si ricostruiscono dal bilancio energetico di superficie,
perché nel bilancio restano tre incognite: i due flussi e il flusso di calore nel suolo, anch'esso non
scritto.

\section{Campi che la configurazione esclude}
\label{sec:configurazione-limiti}

Lo schema YSU è non locale e non porta energia cinetica turbolenta, e con \campo{KM_OPT=4} non la
produce nemmeno la diffusione: il campo non esiste a nessuna quota.

Con \campo{output_diagnostics=0} WRF non ha scritto i massimi e i minimi orari della temperatura a
\SI{2}{m}, e un campionamento orario di \campo{2t} non è l'estremo nell'ora; la simulazione porta
massimi orari, con \campo{nwp_diagnostics=1}, ma non per la temperatura.

Nessun modello d'onda è accoppiato all'atmosfera, e mancano quindi l'altezza significativa, il
periodo e la direzione delle onde, e il coefficiente di resistenza che dipende dallo stato del mare.
L'oceano non è simulato (\campo{SF_OCEAN_PHYSICS=0}): non ci sono correnti né altezza della superficie
del mare. Il ghiaccio marino viene solo dai dati di ingresso, come indicatore a due valori per cella,
con albedo costante \num{0.65}: non esistono lo spessore del ghiaccio né un'albedo del ghiaccio che
evolva.

Con la convezione risolta (\campo{CU_PHYSICS=0}) la precipitazione convettiva \campo{RAINC} è
identicamente nulla, e i parametri convettivi di ERA5, \campo{cp} e \campo{csf}, non si pubblicano
(\cref{sec:non-pubblicati}).

\section{Campi che la regola di pubblicazione esclude}
\label{sec:esclusi}

\paragraph{Componenti della raffica.} Per \campo{10efg} e \campo{10nfg} esiste solo il modulo del
massimo orario del vento risolto (\cref{cap:raffica}), e le due componenti sono escluse per
decisione.

\paragraph{Indice di area fogliare.} Il modello porta un solo indice di area fogliare per cella, un
valore di tabella per categoria dominante (\campo{rdlai2d=.false.}), mentre ERA5 ne vuole uno per la
vegetazione bassa e uno per quella alta. Qualunque divisione del valore unico fra le due sarebbe
un'ipotesi nostra. Il campo \campo{LAI} del modello è nei file wrfout e non è pubblicato.

\paragraph{Angolo e anisotropia dell'orografia sottogriglia.} Il file statico porta i parametri di
asimmetria (\campo{OA1-4}) e di lunghezza efficace (\campo{OL1-4}) dello schema di Kim--Arakawa.
\campo{anor} e \campo{isor} di ECMWF vengono dalla struttura della covarianza del gradiente del
terreno, e le due grandezze sono diverse.

\paragraph{Fusione della neve.} Il campo \campo{ACSNOM} esiste ma non è un accumulatore monotono. In
15 dei 16 giorni esaminati alcuni punti calano fra le 00 e le 12 UTC, fino a \SI{8.2}{mm}, e il 1º
luglio 2024 il calo rispetto alle 00 UTC arriva a \SI{7.26}{mm}. Costruito come gli altri
campi accumulati, come differenza fra due ore successive della stessa run (\cref{sec:tempo}),
\campo{smlt} sarebbe negativo in quei punti.

\section{Campi presenti e non pubblicati}
\label{sec:non-pubblicati}

Sedici campi dei file wrfout sono identicamente nulli in tutti gli istanti del campione di verifica:
\campo{ACHFX}, \campo{ACLHF}, \campo{ACGRDFLX}, \campo{NOAHRES}, \campo{SSTSK}, \campo{SST_INPUT},
\campo{SWNORM}, \campo{ACSNOW}, \campo{HAILNC}, \campo{RAINC}, \campo{RAINSH}, \campo{PREC_ACC_C} e i
quattro campi della radiazione termica discendente al top dell'atmosfera, \campo{LWDNT},
\campo{LWDNTC}, \campo{ACLWDNT} e \campo{ACLWDNTC}. \campo{ACHFX}, \campo{ACLHF}, \campo{ACGRDFLX} e
\campo{NOAHRES} sono nulli perché sono accumulatori dello schema di superficie Noah, e questa
simulazione usa RUC. Un campo identicamente nullo non si pubblica, qualunque sia la ragione.

Per altri campi non nulli nessun parametro ECMWF descrive la stessa grandezza: la frazione congelata
della precipitazione (\campo{SR}), la densità della neve che cade (\campo{RHOSNF}), i massimi e minimi
annuali della frazione di vegetazione verde (\campo{SHDMAX}, \campo{SHDMIN}) e la temperatura del suolo
profondo usata come condizione al contorno inferiore (\campo{TMN}). Il deflusso accumulato
\campo{ACRUNOFF} coincide con il deflusso superficiale \campo{SFROFF}, ed è già in \campo{sro}.

Cinque campi non nulli restano fuori anche se ECMWF definisce un parametro per la stessa grandezza,
perché ERA5 non lo pubblica e l'archivio segue lo schema di ERA5 (\cref{tab:non-pubblicati}).

\begin{table}[htbp]
  \centering
  \small
  \caption{Campi non nulli dei file wrfout con un parametro ECMWF che ERA5 non pubblica.}
  \label{tab:non-pubblicati}
  \begin{tabularx}{\textwidth}{@{}lX@{}}
    \toprule
    Campo WRF & Parametro ECMWF \\
    \midrule
    \campo{SNOALB}     & 260161 \campo{mxsalb}, albedo massima della neve \\
    \campo{QGRAUP}     & 260028 \campo{grle} e 260001 \campo{tcolg}, graupel e sua colonna (\cref{sec:colonne}) \\
    \campo{REFL_10CM}  & 260134 \campo{bref}, riflettività radar \\
    \campo{LPI}        & famiglia 228050 \campo{litoti}; \campo{LPI} è però un indice di potenziale di fulminazione, non una densità di fulmini \\
    \campo{SNOWFALLAC} & 3066 \campo{sde}, perché è uno spessore geometrico della neve e non un'acqua equivalente \\
    \bottomrule
  \end{tabularx}
\end{table}

\chapter{Proprietà che cambiano l'uso dei dati}
\label{cap:conclusioni}

La \cref{tab:proprieta} raccoglie le proprietà dell'archivio di cui tenere conto nel confronto con
ERA5 o con le osservazioni, ciascuna con il valore misurato e la sezione che la descrive.

{\small
\begin{longtable}{@{}p{7.2cm}p{5.4cm}l@{}}
  \caption{Proprietà dell'archivio che cambiano il modo di usarlo.}\label{tab:proprieta}\\
  \toprule
  Proprietà & Valore & Sezione \\
  \midrule
  \endfirsthead
  \toprule
  Proprietà & Valore & Sezione \\
  \midrule
  \endhead
  \bottomrule
  \endfoot
  Radiazione solare discendente in superficie più alta di ERA5; la simulazione non ha aerosol
    & \SI{+8.5}{\percent} in cielo sereno, \SI{+13.6}{\percent} con nubi & \ref{sec:radiazione} \\
  Copertura nuvolosa più bassa di ERA5, per la \campo{CLDFRA} e non per la sovrapposizione
    & rapporto delle medie: \campo{tcc} \num{0.783}, \campo{hcc} \num{0.610} & \ref{sec:colonne} \\
  Temperatura di pelle sulla terra più fredda di ERA5, per cause che l'archivio non separa
    & \SI{-1.64}{K} in media, fino a \SI{-3.24}{K} alle 12 UTC & \ref{sec:pelle} \\
  \campo{fsr} è la rugosità locale di WRF, quella di ERA5 è efficace
    & sulla terra, rapporto delle mediane \numrange{0.32}{0.48}, sei istanti & \ref{sec:rugosita} \\
  \campo{zust} sulla terra risponde al terreno risolto
    & sulla terra, rapporto delle medie \numrange{1.62}{2.06}, sei istanti & \ref{sec:rugosita} \\
  \campo{r} è l'umidità relativa sulla fase mista, \campo{2r} quella sull'acqua liquida
    & fino a un fattore \num{1.8} sotto \SI{250.16}{K} & \ref{sec:livelli-pressione}, \ref{sec:vicino-superficie} \\
  \campo{r} sopra il \SI{100}{\percent} è sovrasaturazione del modello e non è troncata
    & fino al \SI{139.7}{\percent} a \SI{300}{\hPa} il 15 luglio 2024 alle 12 UTC & \ref{sec:livelli-pressione} \\
  L'umidità relativa ricostruita da \campo{2t} e \campo{2d} supera la saturazione
    & fino al \SI{100.09}{\percent} & \ref{sec:vicino-superficie} \\
  Sotto il terreno i livelli di pressione portano il valore del livello del modello più basso (clamping); ERA5 estrapola
    & \SI{1000}{\hPa} sotto il terreno sul \SI{47.6}{\percent} del dominio & \ref{sec:sotto-terreno} \\
  \campo{z} è l'altezza geopotenziale del modello moltiplicata per $g_0 = \SI{9.80665}{m/s^2}$; WRF integra con $g = \SI{9.81}{m/s^2}$
    & \SI{0.034}{\percent} dall'equilibrio idrostatico & \ref{sec:livelli-pressione} \\
  \campo{10fg} è il vento risolto a \SI{10}{m}, non una raffica parametrizzata
    & \num{0.625} volte il \campo{10fg} di ERA5, su 16 istanti & \ref{sec:10fg-era5} \\
  La finestra oraria di \campo{10fg} non è sempre di un'ora
    & azzeramento mancato in 17 dei 138 passaggi orari (\SI{12.3}{\percent}), su sei giorni & \ref{sec:10fg-finestra} \\
  Tre deflussi portano un campo solo; manca il drenaggio a scala giornaliera
    & \campo{ssro} diverso da zero in al più 17 dei \num{2220273} punti, in otto istanti orari del 2024 & \ref{sec:precipitazione} \\
  \campo{tcw} somma sei specie, ECMWF ne definisce cinque
    & residuo di graupel fino a \SI{26.46}{kg/m^2} in quattro istanti del 2024 & \ref{sec:colonne} \\
  Un punto è mancante dove il parametro ERA5 non è definito; bitmap, mai zeri
    & 16 campi con bitmap (\cref{tab:maschere}) & \ref{sec:mascheratura} \\
  \campo{mucin} è mancante dove CAPE è sotto \SI{100}{J/kg}
    & dal \SI{71.2}{\percent} (luglio) al \SI{96.2}{\percent} (gennaio) del dominio, alle 12 UTC & \ref{sec:instabilita} \\
  \campo{generatingProcessIdentifier} non distingue gli accumulati
    & 128 in 246 messaggi su 246 & \ref{sec:tempo} \\
  \campo{al} è un valore di tabella a due stagioni per categoria dominante
    & da 8 a 11 valori sul dominio & \ref{sec:albedo} \\
  \campo{slor} è la pendenza del terreno risolto sotto il nome ECMWF della pendenza sottogriglia; \campo{sdor} è calcolata su celle di \SI{3}{km}
    & \campo{sdor}: rapporto delle mediane \num{0.244} & \ref{sec:orografia} \\
\end{longtable}
}


\appendix
\chapter{Le 90 variabili}
\label{app:variabili}

La \cref{tab:inventario} elenca le 90 variabili di un file orario, lette dai messaggi di un file
orario; ogni file dell'archivio porta gli stessi 246 messaggi (\cref{sec:file}): shortName, paramId e unità sono quelli che eccodes restituisce, e il nome è quello
del parametro ECMWF. Il campionamento è istantaneo, accumulato sull'ora che termina all'istante di validità (\cref{sec:tempo}),
massimo orario (\cref{cap:raffica}) o statico (\cref{sec:statici-orari}). L'ultima colonna indica
i campi con una bitmap di punti mancanti (\cref{sec:mascheratura}); a.e.\ sta per acqua equivalente.

{\footnotesize
\setlength{\tabcolsep}{4pt}
\begin{longtable}{@{}lrp{4.4cm}lp{2.2cm}lc@{}}
  \caption{Le 90 variabili dell'archivio.}\label{tab:inventario}\\
  \toprule
  Campo & paramId & Nome ECMWF & Unità & Livello & Campionamento & Bitmap \\
  \midrule
  \endfirsthead
  \toprule
  Campo & paramId & Nome ECMWF & Unità & Livello & Campionamento & Bitmap \\
  \midrule
  \endhead
  \bottomrule
  \endfoot
  \campo{cc} & 248 & Fraction of cloud cover & 0--1 & 13 livelli di pressione & istantaneo &  \\
  \campo{ciwc} & 247 & Specific cloud ice water content & \si{kg/kg} & 13 livelli di pressione & istantaneo &  \\
  \campo{clwc} & 246 & Specific cloud liquid water content & \si{kg/kg} & 13 livelli di pressione & istantaneo &  \\
  \campo{crwc} & 75 & Specific rain water content & \si{kg/kg} & 13 livelli di pressione & istantaneo &  \\
  \campo{cswc} & 76 & Specific snow water content & \si{kg/kg} & 13 livelli di pressione & istantaneo &  \\
  \campo{q} & 133 & Specific humidity & \si{kg/kg} & 13 livelli di pressione & istantaneo &  \\
  \campo{r} & 157 & Relative humidity & \si{\percent} & 13 livelli di pressione & istantaneo &  \\
  \campo{t} & 130 & Temperature & \si{K} & 13 livelli di pressione & istantaneo &  \\
  \campo{u} & 131 & U component of wind & \si{m/s} & 13 livelli di pressione & istantaneo &  \\
  \campo{v} & 132 & V component of wind & \si{m/s} & 13 livelli di pressione & istantaneo &  \\
  \campo{w} & 135 & Vertical velocity & \si{Pa/s} & 13 livelli di pressione & istantaneo &  \\
  \campo{wz} & 260238 & Geometric vertical velocity & \si{m/s} & 13 livelli di pressione & istantaneo &  \\
  \campo{z} & 129 & Geopotential & \si{m^2/s^2} & 13 livelli di pressione & istantaneo &  \\
  \campo{100u} & 228246 & 100 metre U wind component & \si{m/s} & \SI{100}{m} & istantaneo &  \\
  \campo{100v} & 228247 & 100 metre V wind component & \si{m/s} & \SI{100}{m} & istantaneo &  \\
  \campo{10fg} & 49 & Maximum 10 metre wind gust since previous post-processing & \si{m/s} & \SI{10}{m} & massimo orario &  \\
  \campo{10u} & 165 & 10 metre U wind component & \si{m/s} & \SI{10}{m} & istantaneo &  \\
  \campo{10v} & 166 & 10 metre V wind component & \si{m/s} & \SI{10}{m} & istantaneo &  \\
  \campo{200u} & 228239 & 200 metre U wind component & \si{m/s} & \SI{200}{m} & istantaneo &  \\
  \campo{200v} & 228240 & 200 metre V wind component & \si{m/s} & \SI{200}{m} & istantaneo &  \\
  \campo{2d} & 168 & 2 metre dewpoint temperature & \si{K} & \SI{2}{m} & istantaneo &  \\
  \campo{2r} & 260242 & 2 metre relative humidity & \si{\percent} & \SI{2}{m} & istantaneo &  \\
  \campo{2t} & 167 & 2 metre temperature & \si{K} & \SI{2}{m} & istantaneo &  \\
  \campo{vwsh} & 260068 & Vertical speed shear & \si{s^{-1}} & \SI{100}{m} & istantaneo &  \\
  \campo{msl} & 151 & Mean sea level pressure & \si{Pa} & livello del mare & istantaneo &  \\
  \campo{hcc} & 188 & High cloud cover & 0--1 & colonna & istantaneo &  \\
  \campo{lcc} & 186 & Low cloud cover & 0--1 & colonna & istantaneo &  \\
  \campo{mcc} & 187 & Medium cloud cover & 0--1 & colonna & istantaneo &  \\
  \campo{tcc} & 164 & Total cloud cover & 0--1 & colonna & istantaneo &  \\
  \campo{tciw} & 79 & Total column cloud ice water & \si{kg/m^2} & colonna & istantaneo &  \\
  \campo{tclw} & 78 & Total column cloud liquid water & \si{kg/m^2} & colonna & istantaneo &  \\
  \campo{tcrw} & 228089 & Total column rain water & \si{kg/m^2} & colonna & istantaneo &  \\
  \campo{tcsw} & 228090 & Total column snow water & \si{kg/m^2} & colonna & istantaneo &  \\
  \campo{tcw} & 136 & Total column water & \si{kg/m^2} & colonna & istantaneo &  \\
  \campo{tcwv} & 137 & Total column vertically-integrated water vapour & \si{kg/m^2} & colonna & istantaneo &  \\
  \campo{al} & 174 & Albedo (climatological) & 0--1 & superficie & istantaneo &  \\
  \campo{ci} & 31 & Sea ice area fraction & 0--1 & superficie & istantaneo & sì \\
  \campo{cl} & 26 & Lake cover & 0--1 & superficie & statico &  \\
  \campo{cvh} & 28 & High vegetation cover & 0--1 & superficie & statico &  \\
  \campo{cvl} & 27 & Low vegetation cover & 0--1 & superficie & statico &  \\
  \campo{fal} & 243 & Forecast albedo & 0--1 & superficie & istantaneo & sì \\
  \campo{fsr} & 244 & Forecast surface roughness & \si{m} & superficie & istantaneo &  \\
  \campo{iews} & 229 & Instantaneous eastward turbulent surface stress & \si{N/m^2} & superficie & istantaneo &  \\
  \campo{inss} & 230 & Instantaneous northward turbulent surface stress & \si{N/m^2} & superficie & istantaneo &  \\
  \campo{lsm} & 172 & Land-sea mask & 0--1 & superficie & istantaneo &  \\
  \campo{ro} & 205 & Runoff & \si{m} & superficie & accumulato &  \\
  \campo{rsn} & 33 & Snow density & \si{kg/m^3} & superficie & istantaneo &  \\
  \campo{sd} & 141 & Snow depth & \si{m} a.e. & superficie & istantaneo &  \\
  \campo{sdor} & 160 & Standard deviation of sub-gridscale orography & \si{m} & superficie & istantaneo &  \\
  \campo{sf} & 144 & Snowfall & \si{m} a.e. & superficie & accumulato &  \\
  \campo{skt} & 235 & Skin temperature & \si{K} & superficie & istantaneo &  \\
  \campo{slor} & 163 & Slope of sub-gridscale orography & -- & superficie & istantaneo &  \\
  \campo{slt} & 43 & Soil type & codice & superficie & statico & sì \\
  \campo{snowc} & 260038 & Snow cover & \si{\percent} & superficie & istantaneo &  \\
  \campo{sp} & 134 & Surface pressure & \si{Pa} & superficie & istantaneo &  \\
  \campo{src} & 198 & Skin reservoir content & \si{m} a.e. & superficie & istantaneo &  \\
  \campo{sro} & 8 & Surface runoff & \si{m} & superficie & accumulato &  \\
  \campo{ssr} & 176 & Surface net short-wave (solar) radiation & \si{J/m^2} & superficie & accumulato &  \\
  \campo{ssrc} & 210 & Surface net short-wave (solar) radiation, clear sky & \si{J/m^2} & superficie & accumulato &  \\
  \campo{ssrd} & 169 & Surface short-wave (solar) radiation downwards & \si{J/m^2} & superficie & accumulato &  \\
  \campo{ssrdc} & 228129 & Surface short-wave (solar) radiation downward clear-sky & \si{J/m^2} & superficie & accumulato &  \\
  \campo{ssro} & 9 & Sub-surface runoff & \si{m} & superficie & accumulato &  \\
  \campo{sst} & 34 & Sea surface temperature & \si{K} & superficie & istantaneo & sì \\
  \campo{stl1} & 139 & Soil temperature level 1 & \si{K} & superficie & istantaneo & sì \\
  \campo{stl2} & 170 & Soil temperature level 2 & \si{K} & superficie & istantaneo & sì \\
  \campo{stl3} & 183 & Soil temperature level 3 & \si{K} & superficie & istantaneo & sì \\
  \campo{stl4} & 236 & Soil temperature level 4 & \si{K} & superficie & istantaneo & sì \\
  \campo{str} & 177 & Surface net long-wave (thermal) radiation & \si{J/m^2} & superficie & accumulato &  \\
  \campo{strc} & 211 & Surface net long-wave (thermal) radiation, clear sky & \si{J/m^2} & superficie & accumulato &  \\
  \campo{strd} & 175 & Surface long-wave (thermal) radiation downwards & \si{J/m^2} & superficie & accumulato &  \\
  \campo{strdc} & 228130 & Surface long-wave (thermal) radiation downward clear-sky & \si{J/m^2} & superficie & accumulato &  \\
  \campo{swvl1} & 39 & Volumetric soil water layer 1 & \si{m^3/m^3} & superficie & istantaneo & sì \\
  \campo{swvl2} & 40 & Volumetric soil water layer 2 & \si{m^3/m^3} & superficie & istantaneo & sì \\
  \campo{swvl3} & 41 & Volumetric soil water layer 3 & \si{m^3/m^3} & superficie & istantaneo & sì \\
  \campo{swvl4} & 42 & Volumetric soil water layer 4 & \si{m^3/m^3} & superficie & istantaneo & sì \\
  \campo{tirf} & 235015 & Time integral of rain flux & \si{kg/m^2} & superficie & accumulato &  \\
  \campo{tp} & 228 & Total precipitation & \si{m} & superficie & accumulato &  \\
  \campo{tsn} & 238 & Temperature of snow layer & \si{K} & superficie & istantaneo & sì \\
  \campo{tvh} & 30 & Type of high vegetation & codice 4.234 & superficie & statico &  \\
  \campo{tvl} & 29 & Type of low vegetation & codice 4.234 & superficie & statico & sì \\
  \campo{z} & 129 & Geopotential & \si{m^2/s^2} & superficie & istantaneo &  \\
  \campo{zust} & 228003 & Friction velocity & \si{m/s} & superficie & istantaneo &  \\
  \campo{dl} & 228007 & Lake total depth & \si{m} & lago & statico & sì \\
  \campo{tisr} & 212 & TOA incident short-wave (solar) radiation & \si{J/m^2} & top dell'atmosfera & accumulato &  \\
  \campo{tsr} & 178 & Top net short-wave (solar) radiation & \si{J/m^2} & top dell'atmosfera & accumulato &  \\
  \campo{tsrc} & 208 & Top net short-wave (solar) radiation, clear sky & \si{J/m^2} & top dell'atmosfera & accumulato &  \\
  \campo{ttr} & 179 & Top net long-wave (thermal) radiation & \si{J/m^2} & top dell'atmosfera & accumulato &  \\
  \campo{ttrc} & 209 & Top net long-wave (thermal) radiation, clear sky & \si{J/m^2} & top dell'atmosfera & accumulato &  \\
  \campo{mucape} & 228235 & Most-unstable CAPE & \si{J/kg} & particella più instabile & istantaneo &  \\
  \campo{mucin} & 228236 & Most-unstable CIN & \si{J/kg} & particella più instabile & istantaneo & sì \\
\end{longtable}
}

\chapter{Metodo di verifica}
\label{app:verifica}

Ogni file orario viene confrontato con l'elenco dei 246 messaggi prima di entrare nell'archivio
(\cref{sec:file}). Le proprietà dei capitoli da \ref{cap:atmosfera} a \ref{cap:suolo} sono state
misurate su un campione di verifica di 34 istanti di validità: 4 del 2019, 26 del 2024, che coprono
tutti i mesi, e 4 del 2025. Di questi, 16 sono alle 00 UTC, 16 alle 12 UTC, uno alle 06 e uno alle
18 UTC. Dove un numero è misurato su un insieme diverso, il testo lo dice.\todo{I 16 istanti delle 00 UTC sono stati misurati sulla sesta ora della run del giorno stesso; l'archivio ora le prende dalla trentesima della run del giorno prima (issue 72): da rimisurare}

I confronti con ERA5 usano i campi CHAPTER aggregati sulla griglia di ERA5 con una media pesata
sull'area. Un rapporto fra CHAPTER ed ERA5 è il rapporto delle medie, cioè la media del campo
CHAPTER divisa per la media del campo ERA5 sulla stessa popolazione di punti, oppure, dove indicato,
il rapporto delle mediane.


\printbibliography[heading=bibintoc]

\end{document}
